\documentclass[11pt,reqno,letterpaper]{amsart}
\usepackage[margin=0.9in]{geometry}
\usepackage{amsmath,amssymb,amsthm,mathtools}
\usepackage[T1]{fontenc}
\usepackage[utf8]{inputenc}
\usepackage{lmodern}
\usepackage{microtype,booktabs,longtable,array,enumitem}
\usepackage{needspace}
\usepackage[hidelinks]{hyperref}
\usepackage{bookmark,xurl}
\newcommand{\Z}{\mathbb Z}
\newcommand{\Q}{\mathbb Q}
\newcommand{\C}{\mathbb C}
\newcommand{\PP}{\mathbb P}
\newcommand{\td}{\operatorname{td}}
\newcommand{\ch}{\operatorname{ch}}
\newcommand{\sgn}{\operatorname{sign}}
\newcommand{\Pic}{\operatorname{Pic}}
\newcommand{\Res}{\operatorname{Res}}
\newcommand{\Spec}{\operatorname{Spec}}
\newcommand{\Proj}{\operatorname{Proj}}
\newcommand{\Ahat}{\widehat A}
\newcommand{\OO}{\mathcal O}
\newcommand{\kk}{\mathbf k}
\newtheorem{theorem}{Theorem}[section]
\newtheorem{lemma}[theorem]{Lemma}
\newtheorem{proposition}[theorem]{Proposition}
\newtheorem{conjecture}[theorem]{Conjecture}
\theoremstyle{definition}
\newtheorem{remark}[theorem]{Remark}
\numberwithin{equation}{section}
\allowdisplaybreaks[2]
\setlength{\emergencystretch}{2em}
\setlist{itemsep=2pt,topsep=4pt}
\title[Fano manifolds with the cohomology ring of projective space]{On Fano manifolds with the integral cohomology ring of projective space}
\date{}
\begin{document}
\begin{abstract}
Fujita conjectured that a complex Fano manifold with the integral cohomology
ring of complex projective space is biholomorphic to projective space.
We prove the conjecture in dimensions at most ten. We also prove it in
arbitrary dimension $n$ when the Fano index is at least $n-3$, and in even
dimension when the index is at least $n-7$. The proofs combine
Hirzebruch--Riemann--Roch with bounds on spaces of sections. These give
finitely many possible Hilbert polynomials, which are excluded by
characteristic-number identities and congruences. Two cases require an
additional argument using section rings of genus-two curves. The finite
calculations are carried out by an accompanying Python program.
\end{abstract}
\maketitle

\section{Introduction}\label{sec:intro}

A compact K\"ahler manifold homeomorphic to $\PP^n$ is biholomorphic to
$\PP^n$. This characterization follows from work of Hirzebruch and
Kodaira, Novikov, and Yau; see \cite[Theorem 1]{debarre}. Its proof uses
the rational Pontryagin classes, which are determined by the
homeomorphism type, in a Riemann--Roch calculation. An isomorphism of
integral cohomology rings does not prescribe these classes.

Recall that a \emph{Fano manifold} is a smooth projective complex
manifold whose anticanonical line bundle $K_X^{-1}$ is ample.
Fujita formulated the following conjecture in \cite[\S1]{fujita}.

\begin{conjecture}[Fujita, $C_n$]\label{conj:fujita}
Let $X$ be a Fano manifold of dimension $n$ with
\[
 H^*(X;\Z)\simeq H^*(\PP^n;\Z)
\]
as graded rings. Then $X$ is biholomorphic to $\PP^n$.
\end{conjecture}

The ring structure is essential. An odd-dimensional smooth quadric has
the same integral cohomology groups as projective space, but in dimension
at least three the top power of its ample generator has degree two.
For projective space this degree is one, and it is this restriction that
will give bounds on spaces of sections in our argument.

Fujita proved $C_n$ for $n\le5$ \cite[Theorem 1]{fujita}.
The work of Libgober and Wood \cite{lw} extends the relevant
characteristic-number arguments to dimension six. Debarre
\cite[Theorem 2]{debarre} gives the resulting classification under the
integral cohomology-ring hypothesis: in dimensions at most six, a compact
K\"ahler manifold with this ring is projective space, with the possible
exception of ball quotients in dimensions four and six. The latter have
ample canonical bundle and hence do not occur in the Fano case.
In dimension seven, Debarre's calculation leaves the possibilities
$c_1(X)^7\in\{\pm2^7,\pm4^7\}$ besides projective space
\cite[Theorem 10]{debarre}. For a Fano manifold these are the two positive
indices $2$ and $4$.

Fujita's conjecture arose together with his conjectures $A_n$ and $B_n$
on compactifications and hyperplane sections. More recently, Li and
Peternell \cite{lp} proved that a compact K\"ahler manifold containing a
smooth connected divisor with homologically trivial complement is
projective space with a hyperplane, provided $n\not\equiv3\pmod4$.
Their hypothesis includes a distinguished divisor. Conjecture $C_n$
assumes only the cohomology ring and positivity of the anticanonical
bundle, so such a divisor is not available at the outset.

We prove the following case of Fujita's conjecture.

\begin{theorem}\label{thm:dimension}
Let $X$ be a Fano manifold of dimension $n\le10$ with
$H^*(X;\Z)\simeq H^*(\PP^n;\Z)$ as graded rings.
Then $X$ is biholomorphic to $\PP^n$.
\end{theorem}

\Needspace{8\baselineskip}
The \emph{Fano index} of a Fano manifold $X$ is the
largest positive integer $r$ for which
\[
 K_X^{-1}\simeq A^{\otimes r}
\]
for some holomorphic line bundle $A$. Such a line bundle $A$ is
necessarily ample. Under the cohomology-ring hypothesis of Conjecture
\ref{conj:fujita}, the first Chern class identifies $\Pic(X)$ with
$H^2(X;\Z)\simeq\Z$. If $L$ is its ample generator and $x=c_1(L)$,
then the index is the positive integer $r$ determined by
\[
 K_X^{-1}\simeq L^{\otimes r},\qquad c_1(T_X)=rx,
 \qquad \int_Xx^n=L^n=1.
\]
We also obtain the following result in arbitrary dimension.

\begin{theorem}\label{thm:index}
Let $X$ be a Fano manifold of dimension $n$ with
$H^*(X;\Z)\simeq H^*(\PP^n;\Z)$ as graded rings, and let $r$ be its
Fano index. Then $X$ is biholomorphic to $\PP^n$ if either of the following
conditions holds:
\begin{enumerate}[label=\textup{(\roman*)}]
\item $r\ge n-3$;
\item $n$ is even and $r\ge n-7$.
\end{enumerate}
\end{theorem}

\subsection{Outline of the proof}

By the theorem of Kobayashi and Ochiai \cite{ko}, a Fano $n$-fold has
index at most $n+1$, with equality precisely for $\PP^n$. To prove
Theorems \ref{thm:dimension} and \ref{thm:index}, it therefore suffices
to exclude the possible indices $r<n+1$ under their respective
hypotheses.

We first use the cohomology ring to restrict $r$. The Hodge numbers are
those of $\PP^n$, so $\chi(X,\Omega_X^p)=(-1)^p$. Writing
$c_j(T_X)=c_jx^j$, Riemann--Roch and integrality give
\begin{equation}\label{eq:introindices}
 c_n=n+1,\qquad
 r c_{n-1}=\frac{n(n+1)^2}{2},\qquad
 r\equiv n+1\pmod2.
\end{equation}
In particular, $r$ is a positive divisor of $n(n+1)^2/2$, of the
specified parity, and at most $n+1$. In dimensions seven through ten
the possibilities are as follows.
\[
\begin{array}{c|c|c}
n&\text{indices to exclude}&\text{index of projective space}\\\hline
7&2,4&8\\
8&1,3&9\\
9&2,6&10\\
10&1,5&11
\end{array}
\]

For each remaining index we consider the Hilbert polynomial
$P(t)=\chi(X,L^t)$. Kodaira vanishing gives
\[
 P(-1)=\cdots=P(-(r-1))=0,
 \qquad P(j)=h^0(X,L^j)\quad(j\ge0),
\]
and Serre duality gives $P(-r-t)=(-1)^nP(t)$. These identities, the
leading coefficient $1/n!$, and $P(0)=1$ determine $P$ from the first
few section dimensions $h^0(X,L^j)$; the precise parametrization is
given in Section \ref{sec:hilbert}. The intersection-theoretic bound
\[
 h^0(X,L^j)\le j^n+n
\]
then puts these parameters in finite intervals. No assumption that $L$
is globally generated is required for this bound.

The resulting Hilbert polynomials must satisfy further
characteristic-number identities. In odd dimensions we use
$\chi(X,\Omega_X^p)=(-1)^p$ to eliminate Chern coefficients or Chern
power sums. This produces a polynomial equation in an integer variable,
with coefficients depending on the bounded section dimensions. In even
dimensions we use the signature. The middle intersection form is $(1)$
in the complex orientation, so $\sgn(X)=1$. The Hilbert polynomial
determines the Todd class and hence, after removing the factor
$e^{c_1/2}$, the $\widehat A$-class. A characteristic power-series
identity then recovers the signature and gives a polynomial equation in
the section dimensions.

We exclude integer solutions of these equations by congruences and
discriminant calculations. For $n=7$ and $r=4$, reduction modulo nine
gives a contradiction. For $n=7$ and $r=2$, the discriminant is not a
square for any of the $1224$ allowed pairs of section dimensions.
In the other computational cases, we fix the bounded parameters and
reduce the polynomial in the remaining variable modulo small primes or
prime powers. For each parameter tuple, one of these reductions has
no root. This excludes an integer root without requiring a bound on
the remaining variable.

The assertions in arbitrary dimension also reduce to finitely many
cases. Put $c=n+1-r$. Equation \eqref{eq:introindices} implies
\[
 r\mid\frac{c^2(c-1)}2.
\]
For each fixed positive $c$, only finitely many pairs $(n,r)$ remain.
We treat $c=2,4$, and also $c=6,8$ when $n$ is even. Most of these pairs
are excluded by the same Riemann--Roch or signature calculations. For
$(n,r)=(15,12),(27,24)$, integrality of the twelfth Chern power sum,
using the numerator $691$ of $B_{12}$, forces $h^0(L)=n$ or $n-1$ after
the projective-space case is removed. General sections then cut out an
integral Gorenstein curve of arithmetic genus two. Its section ring
determines a weighted complete-intersection presentation of $X$.
The Euler characteristic of either presentation differs from $n+1$,
giving the required contradiction.

The finite calculations are performed by the accompanying Python file
\nolinkurl{verify_fujita_further_extensions.py}, which we call the
\emph{verification program}. It derives the polynomial equations from
the recurrences in the article and carries out the integer-square and
congruence tests, using exact integer and rational arithmetic.
Section \ref{sec:python} explains the implementation and its relation
to the proofs. The source is distributed separately.

Sections \ref{sec:basic}--\ref{sec:hilbert} establish the numerical and
characteristic-class formulas. Sections \ref{sec:C7}--\ref{sec:C10}
treat dimensions seven through ten. The index bounds are proved in
Sections \ref{sec:high} and \ref{sec:even}. The appendices contain the
longer polynomial expressions and the complete sieve tables.
In dimension eleven, the restrictions above leave $r=2,4,6,8,12$;
our argument excludes $8$, while $12$ gives projective space.
The cases $r=2,4,6$ remain open in this approach.

Throughout, $\PP^n$ denotes complex projective space. We use additive
notation for tensor products of line bundles when convenient, so that
$-K_X=rL$ means $K_X^{-1}\simeq L^{\otimes r}$. Characteristic classes
and their evaluations are taken in the complex orientation.

\section{Numerical restrictions}\label{sec:basic}

\subsection{The ample generator}
Because $X$ is projective, it is K\"ahler. An ample divisor class and each of its
powers are nonzero by positivity. Since $H^{2p}(X;\C)$ has dimension one, that
class spans its $(p,p)$ part. Hodge decomposition therefore gives
\begin{equation}\label{eq:hodge}
 h^{p,q}(X)=
 \begin{cases}1&p=q,\ 0\le p\le n,\\0&\text{otherwise}.\end{cases}
\end{equation}
In particular $H^1(X,\OO_X)=H^2(X,\OO_X)=0$. The exponential sequence implies
\[
 c_1:\Pic(X)\xrightarrow{\ \simeq\ }H^2(X;\Z)\simeq\Z.
\]
One of the two primitive generators corresponds to a line bundle $L$ with an
ample positive power, and hence $L$ itself is ample. Set $x=c_1(L)$.
The ring assumption says that $\int_Xx^n$ is either $1$ or $-1$; positivity of
an ample class in the complex orientation gives
\begin{equation}\label{eq:degreeone}
 \int_Xx^n=1.
\end{equation}
All characteristic classes will be expressed using this choice. Write
$c_j(T_X)=c_jx^j$, so that the symbols $c_j$ on the right are integers.
Since $-K_X$ is ample, its coefficient $r=c_1$ is a positive integer.

\subsection{Holomorphic Euler characteristics}
Equation \eqref{eq:hodge} gives, for every $p$,
\begin{equation}\label{eq:chiOmega}
 \chi(X,\Omega_X^p)=\sum_q(-1)^q h^{p,q}(X)=(-1)^p.
\end{equation}
It also gives $\chi_{\mathrm{top}}(X)=n+1$. The top Chern class is the Euler
class of the underlying real tangent bundle, so
\begin{equation}\label{eq:cn}
 c_n=n+1.
\end{equation}

We derive the relation between $c_1c_{n-1}$ and $c_n$ from the $\chi_y$-genus.
Let $\alpha_1,\ldots,\alpha_n$ be formal Chern roots, and put
$\chi_y(X)=\sum_p\chi(X,\Omega_X^p)y^p$. Riemann--Roch gives
\[
 \chi_{-1+t}(X)=
 \int_X\prod_{i=1}^n
 \alpha_i\left(1+\frac{t}{e^{\alpha_i}-1}\right).
\]
Use
\[
 \frac1{e^\alpha-1}=\frac1\alpha-\frac12+\frac{\alpha}{12}
                  +\text{terms of degree at least three}.
\]
In the coefficient of $t^2$, only the terms with total root-degree zero
after multiplication by $\prod_i\alpha_i$ contribute to the integral.
They give
\[
 [t^2]\chi_{-1+t}(X)=
 \frac14\binom n2 c_n+
 \frac1{12}\left(c_1c_{n-1}-nc_n\right)
 =\frac{n(3n-5)c_n}{24}+\frac{c_1c_{n-1}}{12}.
\]
To see the middle equality, expand
$c_1c_{n-1}=(\sum_j\alpha_j)(\sum_i\prod_{k\ne i}\alpha_k)$;
the $i=j$ terms give $nc_n$.
On the other hand, \eqref{eq:chiOmega} gives
\[
 \chi_{-1+t}(X)=\sum_{p=0}^n(1-t)^p,\qquad
 [t^2]\chi_{-1+t}(X)=\sum_{p=0}^n\binom p2=\binom{n+1}{3}.
\]
Substitute $c_n=n+1$ and solve:
\begin{equation}\label{eq:rcn}
 r c_{n-1}=\frac{n(n+1)^2}{2}.
\end{equation}
In particular, $r$ divides the integer on the right.

\subsection{The Hilbert polynomial and the index}
Let $P(t)=\chi(X,L^t)$. Riemann--Roch gives a polynomial of degree $n$:
\begin{equation}\label{eq:P-leading}
 P(t)=\frac{t^n}{n!}+\frac{rt^{n-1}}{2(n-1)!}+\cdots,\qquad P(0)=1.
\end{equation}
The leading coefficient uses \eqref{eq:degreeone}, and the next coefficient
uses $\td_1=c_1/2$.

For an integer $t>-r$,
\[
 L^t=K_X\otimes L^{t+r}.
\]
The second factor is ample, so Kodaira vanishing gives
$H^i(X,L^t)=0$ for $i>0$. If $t<0$, an anti-ample line bundle has no section:
a nonzero section would give an effective divisor with negative intersection
against $L^{n-1}$. It follows that
\begin{equation}\label{eq:kodaira}
 P(t)=h^0(X,L^t)\quad(t\ge0),\qquad
 P(-1)=\cdots=P(-(r-1))=0.
\end{equation}
There are $r-1$ distinct roots of a nonzero polynomial of degree $n$, hence
\begin{equation}\label{eq:rbound}
 1\le r\le n+1.
\end{equation}
Serre duality also gives the polynomial identity
\begin{equation}\label{eq:SerreP}
 P(-r-t)=(-1)^nP(t).
\end{equation}

The parity restriction can be obtained without any assertion about Pontryagin
classes. The polynomial
\[
 P(t)-\binom{t+n}{n}
\]
is integer-valued on integers and has degree at most $n-1$. Its leading
coefficient is $(r-n-1)/(2(n-1)!)$. The $(n-1)$st forward difference of an
integer-valued polynomial of degree at most $n-1$ is its leading coefficient
times $(n-1)!$, and is an integer. Thus
\begin{equation}\label{eq:parity}
 r\equiv n+1\pmod2.
\end{equation}
Here $\Delta f(t)=f(t+1)-f(t)$;
applying it $d$ times to $at^d$ gives $d!a$.

\subsection{The signature}
If $n$ is even, the middle real cohomology of the real $2n$-manifold $X$ is
$H^n(X;\mathbb R)=\mathbb R x^{n/2}$. Its intersection matrix is the
one-by-one matrix
\[
 \left(\int_Xx^{n/2}x^{n/2}\right)=(1).
\]
Consequently
\begin{equation}\label{eq:signatureone}
 \sgn(X)=1.
\end{equation}
Here the fundamental class is taken in the complex orientation, for which
$\int_Xx^n=1$ by ampleness.

\section{Spaces of sections}\label{sec:sections}

\subsection{The degree bound}
We use the elementary inequality
\begin{equation}\label{eq:min-degree}
 \deg Z\ge N-d+1
\end{equation}
for an integral nondegenerate $d$-dimensional variety $Z\subset\PP^N$.
Here nondegenerate means that $Z$ is not contained in a hyperplane.
One proof cuts by general hyperplanes. Such a section remains nondegenerate:
from $H^0(\OO_Z)=\C$ one has $H^1(\mathcal I_Z)=0$, and the restriction
sequence for the ideal shows that a new linear equation on a hyperplane
section would have come from a linear equation on $Z$.
General sections of positive dimension are integral by Bertini.
After $d$ cuts one has a zero-dimensional scheme of length $\deg Z$
spanning $\PP^{N-d}$. The restrictions of its $N-d+1$ independent linear
coordinates are independent sections of a line bundle on a scheme of that
length; hence the length is at least $N-d+1$.

\begin{lemma}\label{lem:sectionbound}
If $A$ is ample on a smooth projective $n$-fold, then
$h^0(X,A)\le A^n+n$.
\end{lemma}
\begin{proof}
The claim is immediate if $h^0(A)\le1$. Otherwise resolve the complete linear
system by $\mu:Y\to X$ and write
\[
 \mu^*A=M+E,
\]
where $M$ is globally generated and $E$ is effective. Both $\mu^*A$ and $M$
are nef. Let $d$ be the dimension of the image
$Z\subset\PP^{h^0(A)-1}$.
Telescoping powers, each term of the difference
$(\mu^*A)^n-(\mu^*A)^{n-d}M^d$ is an intersection of $E$ with nef classes;
therefore it is nonnegative. The projection formula gives
\[
 A^n\ge(\mu^*A)^{n-d}M^d\ge\deg Z.
\]
The middle intersection is $\deg Z$ times a positive integer: the
$\mu^*A$-degree of a general fiber, or the generic degree if $d=n$.
A general fiber meets the locus where $\mu$ is an isomorphism, so its image
has the same dimension and positive degree with respect to $A$.
By \eqref{eq:min-degree}, $\deg Z\ge h^0(A)-d$.
Since $d\le n$, the result follows.
\end{proof}
Applying this to $A=L^k$ gives
\begin{equation}\label{eq:sectionbound}
 0\le h^0(X,L^k)\le k^n+n\qquad(k\ge1).
\end{equation}

\begin{lemma}\label{lem:base}
Suppose $L^n=1$ and $m=h^0(X,L)>0$. The image of $|L|$ is the full
$\PP^{m-1}$. If the base scheme is nonempty, its codimension is $m$.
\end{lemma}
\begin{proof}
With $A=\mu^*L$, the proof above gives
\[
 1=A^n\ge A^{n-d}M^d\ge\deg Z\ge m-d\ge1.
\]
Every inequality is an equality. In particular, $d=m-1$, so the closed
image in $\PP^{m-1}$ is all of that space.

Let the base scheme have codimension $b<m$. Choose $b$ general members of
$|L|$. They intersect properly: before the $b$th cut the base scheme cannot
force a component of smaller codimension, and general sections avoid all
other components at each stage. Their intersection cycle contains a
base-locus component. It also contains a component outside the base locus,
because $b$ general hyperplanes intersect the dense open image in
$\PP^{m-1}$. These are distinct effective components, each with positive
integral $L$-degree. Their total $L$-degree, however, is $L^n=1$.
This is impossible. Thus $b\ge m$. An ideal generated by the $m$ sections
has height at most $m$, which gives $b=m$.
\end{proof}

\begin{lemma}\label{lem:equality}
If $h^0(X,L)=n+1$, then $(X,L)\simeq(\PP^n,\OO_{\PP^n}(1))$.
Also, if $r=n+1$, then $X\simeq\PP^n$.
\end{lemma}
\begin{proof}
The preceding lemma makes the base scheme empty when $m=n+1$.
The resulting map $f:X\to\PP^n$ has $L=f^*\OO(1)$. Ampleness of $L$
precludes a positive-dimensional fiber, since any curve in such a fiber
would have $L$-degree zero. Thus $f$ is finite. Its degree is $L^n=1$.
A finite birational map to a normal variety is an isomorphism.

If $r=n+1$, the $n$ roots in \eqref{eq:kodaira} and the leading coefficient
of $P$ determine
\[
 P(t)=\frac{(t+1)\cdots(t+n)}{n!}.
\]
Then $h^0(L)=P(1)=n+1$ and the first assertion applies.
\end{proof}
The second assertion is the equality case of the Kobayashi--Ochiai
characterization \cite{ko}; the proof above uses the additional hypothesis
$L^n=1$. For $r<n+1$, the first assertion gives $0\le h^0(L)\le n$.
The weaker bound $n+1$ will suffice in dimensions seven and eight.

\section{Characteristic classes and Riemann--Roch}\label{sec:formal}

\subsection{Formal roots and Newton identities}
Let $s_j$ be the coefficient of $x^j$ in the $j$th power sum of the
Chern roots, with $s_0=n$. Thus $s_1=r$. The $s_j$ are symmetric polynomials in the formal Chern roots.
Newton's identities give
\begin{equation}\label{eq:newton}
 c_k=\frac1k\sum_{j=1}^k(-1)^{j-1}s_jc_{k-j},\qquad c_0=1.
\end{equation}
Conversely,
\[
 s_k=\sum_{j=1}^{k-1}(-1)^{j-1}c_js_{k-j}
       +(-1)^{k-1}kc_k.
\]
The latter formula shows inductively that every $s_k$ is an integer.
The first examples are
$s_2=r^2-2c_2$ and $s_3=r^3-3rc_2+3c_3$.

In all series below, the formal variable $x$ records \emph{complex} degree:
$x^j$ represents a cohomology class of real degree $2j$.
Calculations are truncated after $x^n$.

\subsection{The Todd class}
With Bernoulli numbers defined by $z/(e^z-1)=\sum B_jz^j/j!$,
\begin{equation}\label{eq:logTodd}
 \log\td(T_X)=\frac{rx}{2}
 -\sum_{2j\le n}\frac{B_{2j}s_{2j}}{2j(2j)!}x^{2j}.
\end{equation}
Only even Bernoulli numbers are used. For example, through degree eight,
\[
 \log\td=\frac{rx}{2}-\frac{s_2x^2}{24}
 +\frac{s_4x^4}{2880}-\frac{s_6x^6}{181440}
 +\frac{s_8x^8}{9676800}.
\]
These coefficients come from summing
$\log(\alpha/(1-e^{-\alpha}))$ over the formal roots.

We shall use the following recurrences for truncated power series. If
$A(x)=1+\sum_{j\ge1}a_jx^j$ and $\log A(x)=\sum_{j\ge1}b_jx^j$, then
\begin{equation}\label{eq:logrec}
 b_j=a_j-\frac1j\sum_{k=1}^{j-1}k b_k a_{j-k}.
\end{equation}
If $A(x)=\exp B(x)$, the converse is
\begin{equation}\label{eq:exprec}
 a_0=1,\qquad a_j=\frac1j\sum_{k=1}^j k b_k a_{j-k}.
\end{equation}
Both follow from $A'=B'A$. Products are computed by the usual finite
convolution $[x^j](UV)=\sum_{k=0}^ju_kv_{j-k}$.
In particular, all these operations take place over $\Q$.

\subsection{Exterior powers of the cotangent bundle}
If $\td(T_X)=\sum_{j=0}^n\td_jx^j$, Riemann--Roch is
\begin{equation}\label{eq:recoverTodd}
 P(t)=[x^n]e^{tx}\td(T_X),\qquad
 \td_j=(n-j)![t^{n-j}]P(t).
\end{equation}
The Chern character of the cotangent bundle is
\[
 U(x)=\ch(\Omega_X^1)=\sum_{j=0}^n\frac{(-1)^js_j}{j!}x^j.
\]
Writing $U_p(x)=\ch(\Lambda^p\Omega_X^1)$, with $U_0=1$, Newton's
exterior-power identity is
\begin{equation}\label{eq:exterior}
 p U_p(x)=\sum_{a=1}^p(-1)^{a-1}U_{p-a}(x)\,U(ax).
\end{equation}
Indeed, the roots of the exterior-power generating function are
$e^{-\alpha_1},\ldots,e^{-\alpha_n}$; their $a$th power sum is $U(ax)$.
For example, $U_2=(U(x)^2-U(2x))/2$.
Finally,
\begin{equation}\label{eq:HRRexterior}
 [x^n]U_p(x)\td(T_X)=\chi(X,\Omega_X^p)=(-1)^p.
\end{equation}
We will use these recurrences to express the Riemann--Roch identities
as polynomial equations in the Chern power sums.

\section{The Hilbert series and the signature}\label{sec:hilbert}

\subsection{Reciprocity of the Hilbert numerator}
Put $H(z)=\sum_{k\ge0}P(k)z^k$. Finite differences of a degree-$n$
polynomial show that
\[
 H(z)=\frac{h(z)}{(1-z)^{n+1}},\qquad h(z)\in\Z[z],\quad\deg h\le n.
\]
For example, the coefficient of $z^j$, $j\ge n+1$, in
$(1-z)^{n+1}H(z)$ is the $(n+1)$st difference of $P$, and is zero.
We have $h(0)=P(0)=1$. The leading coefficient $1/n!$ of $P$ gives
$h(1)=1$ by comparison of the leading pole at $z=1$.

The rational generating function of a polynomial satisfies
\[
 H(z^{-1})=-\sum_{k\ge1}P(-k)z^k.
\]
The equality means expansion of the rational function on the left at zero.
It follows first for $P(t)=1$ from
$(1-z^{-1})^{-1}=-z/(1-z)$; applying $z\,d/dz$ proves it for monomials,
and hence for every polynomial. Using \eqref{eq:kodaira} and
\eqref{eq:SerreP} on the right gives
\[
 H(z^{-1})=(-1)^{n+1}z^rH(z).
\]
Substituting the rational form of $H$ and cancelling the denominators yields
\begin{equation}\label{eq:reciprocal}
 h(z)=z^{n+1-r}h(z^{-1}).
\end{equation}
Since $h(0)=1$, the degree is exactly $n+1-r$, its leading coefficient is
one, and its coefficients are symmetric.
By \eqref{eq:parity}, write $n+1-r=2s$.

Let $v_i=h^0(X,L^i)$, with $v_0=1$, and $h(z)=\sum_{j=0}^{2s}h_jz^j$.
Coefficient comparison gives
\begin{align}
 h_i&=v_i-\sum_{j=0}^{i-1}h_j\binom{n+i-j}{n}\quad(0\le i<s),\label{eq:hrec}\\
 h_s&=1-2\sum_{j=0}^{s-1}h_j,\qquad h_{2s-j}=h_j.\nonumber
\end{align}
Thus only $v_1,\ldots,v_{s-1}$ are free parameters.
We make no assertion that the coefficients $h_j$ are nonnegative:
$L$ need not generate its section ring in degree one.

\subsection{Recovering the Todd class}
The following identity holds modulo $x^{n+1}$:
\begin{equation}\label{eq:Todd-h}
 \td(T_X)=\left(\frac{x}{1-e^{-x}}\right)^{n+1}h(e^{-x}).
\end{equation}
The binomial coefficient identity follows by a formal residue calculation.
For any integer $j$,
\[
 [x^n]e^{(t-j)x}\left(\frac{x}{1-e^{-x}}\right)^{n+1}
 =\Res_{x=0}\frac{e^{(t-j)x}\,dx}{(1-e^{-x})^{n+1}}.
\]
Put $z=1-e^{-x}$, so $dx=dz/(1-z)$. The residue becomes
$[z^n](1-z)^{-(t-j+1)}=\binom{t+n-j}{n}$.
Multiplication by $h_j$ and summation gives exactly the Hilbert polynomial.
Since \eqref{eq:recoverTodd} recovers each coefficient separately, the
truncated classes are equal.

\subsection{The signature formula}
Set
\[
 q(x)=\frac{x}{2\sinh(x/2)},\qquad B(x)=e^{sx}h(e^{-x}).
\]
Reciprocity makes $B$ even, and \eqref{eq:Todd-h} gives
\[
 \Ahat(T_X)=e^{-rx/2}\td(T_X)=q(x)^{n+1}B(x).
\]
This is an identity of rational characteristic classes; it requires no
spin structure. The characteristic power series obey
\[
 \frac{q(2x)^2}{q(4x)}=\frac{x}{\tanh x}.
\]
Consequently, root by root and then by multiplication, the Hirzebruch
$L$-class is
\[
 \mathcal L(T_X)=\left(\frac{x}{\tanh x}\right)^{n+1}
                 \frac{B(2x)^2}{B(4x)}.
\]
The factors $e^{2sx}$ and $e^{4sx}$ cancel in this quotient. Taking the
top coefficient and making the formal change of variable $t=\tanh x$
gives
\begin{equation}\label{eq:signature-h}
 \sgn(X)=[t^n]\frac{
 h\!\left(\frac{1-t}{1+t}\right)^2}{
 (1-t^2)h\!\left(\left(\frac{1-t}{1+t}\right)^2\right)}.
\end{equation}
Indeed, the coefficient is initially the residue of
$h(e^{-2x})^2\,dx/(\tanh(x)^{n+1}h(e^{-4x}))$.
Now $dx=dt/(1-t^2)$ and $e^{-2x}=(1-t)/(1+t)$.

To express this in terms of polynomials, put $d=2s$ and define
\begin{align}
 N(t)&=\sum_{j=0}^d h_j(1-t)^j(1+t)^{d-j},\label{eq:ND}\\
 D(t)&=(1-t^2)\sum_{j=0}^d h_j(1-t)^{2j}(1+t)^{2d-2j}.\nonumber
\end{align}
Both are even, and $N(0)=D(0)=1$. Thus, in even complex dimension,
\begin{equation}\label{eq:sigmaND}
 [t^n]\frac{N(t)^2}{D(t)}=1.
\end{equation}
If $D(t)=1+\sum_{j\ge1}d_jt^{2j}$ and
$N(t)^2=\sum_{j\ge0}b_jt^{2j}$, set
\begin{equation}\label{eq:division}
 a_0=b_0=1,\qquad a_j=b_j-\sum_{i=1}^j d_i a_{j-i}.
\end{equation}
Then \eqref{eq:sigmaND} is $a_{n/2}=1$. Every coefficient in this recurrence
is an integer polynomial in the section counts.

\subsection{Reduction modulo an integer}
Let $F\in\Z[y]$. If $F$ has an integer root, its reduction modulo
any positive integer $q$ has a root in $\Z/q\Z$. We shall use the
contrapositive: a modulus for which the reduction has no root excludes
all integer roots of $F$. When the coefficients depend on a bounded
integer tuple $\boldsymbol v$, we may choose a different modulus for
each specialization $F_{\boldsymbol v}$.

We first clear denominators and divide out the greatest common divisor
of the coefficients over $\Q$, then reduce the resulting integer
polynomial modulo $q$. These operations preserve its rational zero set.
In particular, a prime dividing a denominator in the derivation may
still be used to exclude integer roots of the resulting polynomial.

\section{Dimension seven}\label{sec:C7}

\subsection{The possible indices}
In this section only, write the Chern coefficients as
\[
 (c_1,c_2,c_3,c_4,c_5,c_6,c_7)=(r,a,b,c,d,e,f).
\]
In particular, $b=c_3$; the power sum $s_3$ retains the notation of
Section \ref{sec:formal}.
Equations \eqref{eq:cn}--\eqref{eq:parity} become
\[
 f=8,\qquad re=224,\qquad 1\le r\le8,\qquad r\text{ even}.
\]
Among $2,4,6,8$, the value $6$ does not divide $224$. Hence
\begin{equation}\label{eq:C7indices}
 r\in\{2,4,8\}.
\end{equation}
The value $8$ gives $\PP^7$ by Lemma \ref{lem:equality}.
It remains to exclude $4$ and $2$.

\subsection{The Chern-number equations}
Apply the recurrences of Section \ref{sec:formal}, substituting $e=224/r$
and $f=8$. Let
\[
 \rho_0=\chi(\OO_X)-1,\quad
 \rho_1=\chi(\Omega_X^1)+1,\quad
 \rho_2=\chi(\Omega_X^2)-1.
\]
The resulting identities are
\[
 E_1=-7200\rho_0+1440\rho_1,\qquad
 E_2=-120960\rho_0,\qquad
 \rho_2=-6\rho_0+3\rho_1,
\]
where
\begin{align}
 E_1={}&r^3c-3rac-r^2d+3bc-3ad+7728,\label{eq:C7E1}\\
 E_2={}&-2r^5a+10r^3a^2+2r^4b-10ra^3-11r^2ab\nonumber\\
       &-2r^3c+rb^2+9rac+2r^2d+120512.\label{eq:C7E2}
\end{align}
Under the substitutions $e=224/r$ and $f=8$, the three
Euler-characteristic identities are therefore equivalent to
$E_1=E_2=0$.

For use in the elimination, the Hilbert polynomial is
\begin{align}
60480P(t)={}&12t^7+42rt^6+42(r^2+a)t^5+105rat^4\nonumber\\
 &+14(-r^4+4r^2a+3a^2+rb-c)t^3\nonumber\\
 &+21(-r^3a+3ra^2+r^2b-rc)t^2\nonumber\\
 &+\bigl(2r^6-12r^4a+11r^2a^2+5r^3b+10a^3\nonumber\\
 &\qquad+11rab-5r^2c-b^2-9ac-2rd+2e\bigr)t+60480.
 \label{eq:C7P}
\end{align}
The constant term uses $\rho_0=0$. The remaining coefficients follow
from \eqref{eq:logTodd}, \eqref{eq:exprec}, and
\eqref{eq:recoverTodd}, truncated in degree seven. For example, the coefficient
of $t^5$ is $\td_2/5!=(r^2+a)/(12\cdot5!)$,
which is $42(r^2+a)/60480$ as displayed.

\subsection{The case \texorpdfstring{$r=4$}{r=4}}
Set $r=4$, $e=56$, and $m=P(1)$. The coefficient of $d$ in $E_2$
is $2r^2=32$, so solve $E_2=0$ uniquely for $d$.
Substitution in \eqref{eq:C7P} gives
\[
 P(-1)=-\frac{3a^2-41a+4b-c-598}{1440},\qquad
 P(1)=\frac{3a^2-17a+4b-c+458}{288}.
\]
Kodaira vanishing gives $P(-1)=0$. Subtract the corresponding numerators:
$288m=24a+1056$, hence
\[
 a=12m-44,\qquad
 c=4b+432m^2-3660m+7014.
\]
The solved equation for $d$ becomes
\[
 8d=-b^2+96bm-352b-29376m^3
       +367056m^2-1451928m+1787736.
\]
Substituting these expressions in $2E_1=0$ gives
\begin{align}
 F_4(m,b)={}&(9m-5)b^2+(1728m^2-17160m+36612)b\nonumber\\
 &+264384m^4-4279824m^3+25277544m^2\nonumber\\
 &-64549008m+60164376=0.\label{eq:C7F4}
\end{align}
Only nonzero constants have been divided out in these substitutions.

Reduce all coefficients modulo nine:
\[
 F_4(m,b)\equiv4b^2+3bm+6\pmod9.
\]
Reducing further modulo three forces $b^2\equiv0\pmod3$, so $3\mid b$.
Then both $4b^2$ and $3bm$ are divisible by nine, while $6$ is not.
This contradiction excludes $r=4$, for every integer $m$ without any bound.

\subsection{The case \texorpdfstring{$r=2$}{r=2}}
Now $r=2$ and $e=112$. Set $m=P(1)$ and $k=P(2)$.
Equation \eqref{eq:sectionbound} gives
\begin{equation}\label{eq:C7bounds}
 0\le m\le8,\qquad 0\le k\le 2^7+7=135.
\end{equation}
Again solve $E_2=0$ for $d$, this time dividing by the nonzero constant $8$.
The two evaluations of $P$ reduce to
\[
 P(1)=\frac{3a^2+a+2b-c+1442}{720},\qquad
 P(2)=\frac{3a^2+16a+2b-c+572}{180}.
\]
Thus $180k-720m=15a-870$. Put $u=k-4m$; then
\begin{align*}
 a&=12u+58,\\
 c&=2b+432u^2+4188u+11592-720m,\\
 4d&=-b^2+48bu+232b+77760mu+370080m\\
    &\quad-29376u^3-429552u^2-2248824u-4200128.
\end{align*}
With these substitutions, $2E_1=0$ becomes
\begin{equation}\label{eq:C7quad}
 A b^2+B b+C=0,
\end{equation}
where
\begin{align*}
 A&=18u+101,\\
 B&=1728u^2+16392u+47544-4320m,\\
 C&=528768u^4+10284192u^3+77812128u^2\\
  &\quad+271230552u+365944288\\
  &\quad-m(1399680u^2+13478400u+32447520).
\end{align*}
The integer $A$ cannot vanish, since $101$ is not divisible by $18$.
An integral root $b$ would give the identity
\begin{equation}\label{eq:disc}
 B^2-4AC=(2Ab+B)^2.
\end{equation}
Therefore the discriminant must be a nonnegative integer square.

The rectangle \eqref{eq:C7bounds} contains $9\cdot136=1224$ pairs.
Computing $u,A,B,C$ and $D=B^2-4AC$ for each pair gives
\[
\begin{array}{c|rrrrrrrrr|r}
m&0&1&2&3&4&5&6&7&8&\text{total}\\ \hline
D\ge0&0&0&4&7&12&16&20&25&30&114\\
D\text{ square}&0&0&0&0&0&0&0&0&0&0
\end{array}
\]
For each nonnegative discriminant, we compute the integer
$\lfloor\sqrt D\rfloor$ and compare its square with $D$.
Alternatively, requiring the $114$ nonnegative discriminants to be
quadratic residues successively modulo $5,7,11,13,17,19,23,29$ leaves
$90,60,33,23,16,8,3,0$ pairs. Either calculation excludes every pair
in \eqref{eq:C7bounds}.

Both nonmaximal indices in \eqref{eq:C7indices} are excluded.
The remaining value $r=8$ gives $X\simeq\PP^7$ by Lemma \ref{lem:equality}.
This proves $C_7$.

\section{Dimension eight}\label{sec:C8}

\subsection{The possible indices}
Now $c_8=9$ and $r c_7=324$. Bounds and parity give the candidates
$1,3,5,7,9$; divisibility by $324$ removes $5,7$. Hence
\begin{equation}\label{eq:C8indices}
 r\in\{1,3,9\}.
\end{equation}
The value $9$ gives $\PP^8$ by Lemma \ref{lem:equality}.

\subsection{The Hilbert polynomial}
Since $n$ is even, \eqref{eq:SerreP} says $P(t)=P(-r-t)$.
Put $y=t+r/2$. Invariance under $y\mapsto-y$ means that $P$ is a polynomial
in $y^2$, equivalently in $z=t(t+r)=y^2-r^2/4$.
The leading coefficient and $P(0)=1$ give
\begin{equation}\label{eq:C8P}
 P(t)=\frac{z^4+A z^3+B z^2+Cz+40320}{40320},
 \qquad z=t(t+r).
\end{equation}
The letters $A,B,C$ in this section denote Hilbert-polynomial coefficients.

Set $m=P(1)$, $k=P(2)$, $\ell=P(3)$. The bounds needed are
\begin{equation}\label{eq:C8bounds}
 0\le m\le9,\qquad 0\le k\le2^8+8=264.
\end{equation}
For $r=3$ the values of $z$ at $t=-1,1,2$ are $-2,4,10$.
The three equations for $A,B,C$ are therefore
\[
\begin{aligned}
 -8A+4B-2C&=-40336,\\
 64A+16B+4C&=40320m-40576,\\
 1000A+100B+10C&=40320k-50320.
\end{aligned}
\]
Their coefficient matrix is invertible, as is also clear from polynomial
interpolation at three distinct nonzero values of $z$. Solving gives
\begin{equation}\label{eq:C8abc3}
\begin{aligned}
 A&=56k-280m+492,\\
 B&=-112k+2240m-6036,\\
 C&=-448k+5600m+6128.
\end{aligned}
\end{equation}
For $r=1$, the values of $z$ at $t=1,2,3$ are $2,6,12$.
The equations are
\[
 8A+4B+2C=40320m-40336,\quad
 216A+36B+6C=40320k-41616,
\]
\[
 1728A+144B+12C=40320\ell-61056.
\]
Their solution is
\begin{equation}\label{eq:C8abc1}
\begin{aligned}
 A&=-280k+56\ell+504m-300,\\
 B&=3920k-448\ell-9072m+5708,\\
 C&=-6720k+672\ell+36288m-30384.
\end{aligned}
\end{equation}

\subsection{The signature equation}
In dimension eight we can write the signature in terms of the four
coefficients of the $\widehat A$-class. By \eqref{eq:recoverTodd},
$\td_j=(8-j)![t^{8-j}]P(t)$. Write
\[
 e^{-rx/2}\sum_{j=0}^8\td_jx^j
 =1+Ux^2+Vx^4+Wx^6+Zx^8.
\]
The odd coefficients vanish, either by \eqref{eq:SerreP} or by the
Todd power-series identity. The logarithm has coefficients
\begin{align*}
 q_1&=U,\\
 q_2&=V-U^2/2,\\
 q_3&=W-UV+U^3/3,\\
 q_4&=Z-UW-V^2/2+U^2V-U^4/4.
\end{align*}
The conversion to the logarithm of the $L$-class is:
\[
\begin{array}{c|c|c|c}
 &\log\Ahat\text{ coefficient of }s_{2j}
 &\log\mathcal L\text{ coefficient of }s_{2j}&\text{ratio}\\ \hline
j=1&-1/24&1/3&-8\\
j=2&1/2880&-7/90&-224\\
j=3&-1/181440&62/2835&-3968\\
j=4&1/9676800&-127/18900&-65024
\end{array}
\]
Each entry follows by expanding $\log(x/(2\sinh(x/2)))$ or
$\log(x/\tanh x)$ and summing over roots.
Set $(\lambda_1,\lambda_2,\lambda_3,\lambda_4)
=(-8q_1,-224q_2,-3968q_3,-65024q_4)$.
Exponentiating and extracting degree eight gives
\begin{equation}\label{eq:C8sigma}
 \sgn(X)=\lambda_4+\lambda_1\lambda_3+\frac{\lambda_2^2}{2}
              +\frac{\lambda_1^2\lambda_2}{2}+\frac{\lambda_1^4}{24}.
\end{equation}
By \eqref{eq:signatureone}, this number must be $1$.

\subsection{The case \texorpdfstring{$r=3$}{r=3}}
Substitute \eqref{eq:C8abc3} into \eqref{eq:C8P} and then into
\eqref{eq:C8sigma}. Set $u=k-5m$. We obtain the polynomial identity
$\sgn(X)-1=32F_3(m,u)$, where
\begin{align}
 F_3(m,u)={}&1152u^4+40920u^3+(557147-4032m)u^2\nonumber\\
 &+(3443306-71232m)u+1800m^2-325410m+8148917.
 \label{eq:C8F3}
\end{align}
Thus $F_3(m,u)=0$ for an integer $u$ and $m=0,\ldots,9$.
For each value of $m$, a modulus excluding every residue of $u$ is:
\[
\begin{array}{c|rrrrrrrrrr}
m&0&1&2&3&4&5&6&7&8&9\\ \hline
q&11&11&7&7&11&7&7&19&11&7
\end{array}
\]
For each entry, $F_3(m,u)$ is nonzero modulo $q$ for
$u=0,\ldots,q-1$. Hence there is no integer root for any of the
allowed values of $m$, and $r=3$ is impossible. The argument uses
no bound on $u$ or $k$.

\subsection{The case \texorpdfstring{$r=1$}{r=1}}
For \eqref{eq:C8abc1}, set $u=\ell-5k+9m$ and $v=k-3m$.
The same calculation now gives $\sgn(X)-1=32F_1(m,k,u)$, where
\begin{align}
 F_1(m,k,u)={}&1152u^4-23592u^3+173099u^2-537930u+596995\nonumber\\
 &+1800v^2+(-4032u^2+41664u-100434)v\nonumber\\
 &+(3024u-16000)m,\qquad v=k-3m.\label{eq:C8F1}
\end{align}
The two bounds in \eqref{eq:C8bounds} give $10\cdot265=2650$ pairs $(m,k)$.
For each pair, treat \eqref{eq:C8F1} as a polynomial in the unrestricted
integer $u$. The successive congruence tests give the following counts:
\[
\begin{array}{r|r@{\qquad}r|r}
q&\text{remaining pairs}&q&\text{remaining pairs}\\ \hline
\text{initial}&2650&13&39\\
4&1325&17&25\\
9&662&19&14\\
16&332&23&8\\
25&160&29&3\\
27&148&31&2\\
49&80&37&1\\
11&64&43&0
\end{array}
\]
The moduli are applied down the left column and then down the right.
As explained in Section \ref{sec:hilbert}, an integer root must also
give a root modulo each composite modulus. No pairs remain, so $r=1$
is impossible for every integer $\ell$.
Equation \eqref{eq:C8indices} leaves $r=9$, and Lemma \ref{lem:equality}
gives $X\simeq\PP^8$. This proves $C_8$.

\section{Dimension nine}\label{sec:C9}
The basic restrictions give $r\mid450$, $r$ even, and $r\le10$.
Thus $r=2,6,10$. The value $10$ gives projective space.
The value $6$ is excluded in Section \ref{sec:higharith} by direct
characteristic-number equations; that argument does not assume $C_9$.
We treat $r=2$ here.

\subsection{Section dimensions and even power sums}
Set $m=h^0(L)$, $k=h^0(2L)$, and $\ell=h^0(3L)$.
Lemma \ref{lem:equality} excludes $m=10$ in the present index, so
\begin{equation}\label{eq:C9bounds}
 0\le m\le9,\quad 0\le k\le 2^9+9=521,\quad
 0\le\ell\le3^9+9=19692.
\end{equation}
Since $n+1-r=8$, the first half of $h(z)$ is
\[
 h_0=1,\quad h_1=m-10,\quad h_2=k-10m+45,\quad
 h_3=\ell-10k+45m-120.
\]
Set $h_4=1-2(h_0+h_1+h_2+h_3)$ and reflect the coefficients.

In this section put
\[
 u=s_2,\quad v=s_4,\quad w=s_6,\quad a=s_3,\quad b=s_5.
\]
Thus the letters $a,b$ in this section denote odd Chern power sums. Let $B(x)=e^{4x}h(e^{-x})$. If
\[
 T=\ell-6k+14m-14,\quad Q=\ell+6k-34m+46,\quad
 J=\ell+54k+134m-434,
\]
direct expansion of the paired exponentials gives
\[
 [x^2]B=T,\qquad [x^4]B=Q/12,\qquad [x^6]B=J/360.
\]
For example, $B=2\cosh4x+2h_1\cosh3x+2h_2\cosh2x
+2h_3\cosh x+h_4$, so its degree-two coefficient is
$16+9h_1+4h_2+h_3=T$.

Since $\log\Ahat=10\log q+\log B$, formula \eqref{eq:logrec} gives
\begin{equation}\label{eq:C9powers}
\begin{aligned}
 u&=10-24T,\\
 v&=10+240Q-1440T^2,\\
 w&=10-504J+15120TQ-60480T^3.
\end{aligned}
\end{equation}
For instance, $[x^4]\log B=Q/12-T^2/2$, and the coefficient of
$s_4$ in $\log q$ is $1/2880$.

\subsection{Elimination of the odd power sums}
Start from
\[
 (s_0,\ldots,s_9)=(9,2,u,a,v,b,w,s_7,s_8,s_9).
\]
The equation $[x^9]\td=1$ is linear in $s_8$ with coefficient
$1/9676800$, hence has the unique solution
\begin{equation}\label{eq:s8nine}
\begin{aligned}
s_8={}&-\frac{175}{144}u^4+\frac{175}{9}u^3-\frac{35}{12}u^2v-70u^2\\
 &+\frac{70}{3}uv-\frac{20}{9}uw+80u-\frac7{12}v^2\\
 &-28v+\frac{80}{9}w+\frac{29030320}{3}.
\end{aligned}
\end{equation}
Indeed, $s_8x^8/9676800$
in $\log\td$ multiplies the degree-one term $x$ of $\td$.

Newton's identities now compute $c_8$ and $c_9$. They must be
\[
 c_8=\frac{9\cdot10^2}{2\cdot2}=225,\qquad c_9=10.
\]
The coefficient of $s_7$ in $c_8$ is $s_1/7=2/7$.
Thus, if $c_8^{(0)}$ denotes the Newton expression with $s_7=0$,
\[
 s_7=\frac72(225-c_8^{(0)}).
\]
After this substitution the coefficient of $s_9$ in $c_9$ is $1/9$,
so, with the analogous notation, $s_9=9(10-c_9^{(0)})$.
These are linear eliminations with nonzero constant coefficients.
They are therefore valid for every choice of $u,v,w,a,b$.

Impose \eqref{eq:HRRexterior} for $p=1,2$. Clearing denominators and
coefficient gcds gives two integer polynomials
\[
 E_1(u,v,w,a,b)=0,\qquad E_2(u,v,w,a,b)=0.
\]
Each has degree one in $b$. Write $E_i=L_i b+C_i$. Their resultant is
$L_1C_2-L_2C_1$. Any common solution makes this expression zero, even if
one or both coefficients $L_i$ vanish. No division by $L_i$ is needed.

Remove the resultant's nonzero coefficient gcd. The resulting polynomial
$F(u,v,w,a)$ has 60 monomials and degree four in $a$:
\begin{equation}\label{eq:C9quartic}
 F(u,v,w,a)=\sum_{i=0}^4 A_i(u,v,w)a^i,\qquad
 A_4=3840u-63488.
\end{equation}
The equations $E_1,E_2$ and the five coefficients $A_i$ are given in
Appendix \ref{app:c9poly}. After substituting \eqref{eq:C9powers},
the integer $a=s_3$ must satisfy \eqref{eq:C9quartic}.

\subsection{The congruence calculation}
The rectangle \eqref{eq:C9bounds} contains
$10\cdot522\cdot19693=102797460$ triples.
For each triple, the integer polynomial in $a$ must have a root modulo
every prime. Test the primes from $13$ through $223$:
\[
\begin{array}{r|r}
 \text{primes tested through}&\text{remaining triples}\\ \hline
23&20583262\\97&16021\\149&136\\181&3\\223&0
\end{array}
\]
Appendix \ref{app:sievetables} contains every intermediate count.
The representation of the candidate set used for this calculation is
described in Section \ref{sec:python}.
The empty final set excludes $r=2$. Only $r=10$ remains, proving $C_9$.

\section{Dimension ten}\label{sec:C10}
Here $r\mid605$, $r$ is odd, and $r\le11$. The three possibilities are
$r=1,5,11$. Section \ref{sec:even6} excludes $r=5$ by a signature
calculation independent of the present argument. The value $11$ gives
projective space. It remains to exclude $r=1$.

\subsection{The signature polynomial}
Put $m=h^0(L)$, $k=h^0(2L)$, $\ell=h^0(3L)$, and $j=h^0(4L)$.
Now $h$ has degree ten, with
\begin{align*}
 h_0&=1,&h_1&=m-11,\\
 h_2&=k-11m+55,&h_3&=\ell-11k+55m-165,\\
 h_4&=j-11\ell+55k-165m+330.
\end{align*}
Its central coefficient and its reflected coefficients are given by
\eqref{eq:hrec}. Formula \eqref{eq:division} gives
\begin{equation}\label{eq:C10F}
 \sgn(X)-1=8F_{10,1}(m,k,\ell,j),
\end{equation}
where $F_{10,1}$ is a primitive integer polynomial with 126 monomials
and degree five in $j$. Its leading coefficient in $j$ is $-73728$.

The polynomial can be expressed more briefly in terms of the following
linear combinations:
\begin{align*}
 T&=j-7\ell+20k-28m+14,\\
 Q&=j+5\ell-40k+80m-46,\\
 J&=j+53\ell+80k-568m+434,\\
 V&=j+245\ell+3800k+7280m-11326.
\end{align*}
They are respectively $[x^2]B$, $12[x^4]B$, $360[x^6]B$,
and $20160[x^8]B$ for $B(x)=e^{5x}h(e^{-x})$.
Let
\begin{align}
G(T,Q,J,V)={}&-23224320T^5+5322240T^4-532224T^3\nonumber\\
 &+30162T^2+12151T-41207040\nonumber\\
 &+8709120QT^3-1552320QT^2+110880QT-22533Q\nonumber\\
 &-655200Q^2T+57750Q^2\nonumber\\
 &-233856JT^2+38808JT+3654J+17640JQ\nonumber\\
 &+3060TV-1147V.\label{eq:C10compact}
\end{align}
Then
\begin{equation}\label{eq:C10division}
 F_{10,1}(m,k,\ell,j)=\frac{G(T,Q,J,V)}{315}.
\end{equation}
After the displayed linear substitutions, all coefficients of $G$ are
divisible by $315$. The resulting integer polynomial agrees with that
obtained from \eqref{eq:division}, and also with the expression
recovered from the Todd class. Both identities are checked by the
verification program.

The congruence tests are applied to $F_{10,1}$. The division by $315$
is performed coefficientwise over $\Q$ before reduction modulo a prime,
as in Section \ref{sec:hilbert}.

\subsection{The congruence calculation}
The bounds required for the first three section counts are
\[
 0\le m\le10,\qquad0\le k\le2^{10}+10=1034,\qquad
 0\le\ell\le3^{10}+10=59059.
\]
The value $m=11$ has been removed by Lemma \ref{lem:equality}.
No bound on $j$ is needed, since an integer root gives a root modulo
every prime.

There are $11\cdot1035\cdot59060=672398100$ triples. The congruence tests give
\[
\begin{array}{r|r}
 \text{primes tested through}&\text{remaining triples}\\ \hline
31&8661101\\97&14926\\149&148\\181&7\\239&0.
\end{array}
\]
Thus $F_{10,1}$ has no integer root in $j$ for any allowed triple
$(m,k,\ell)$.
This excludes $r=1$, and only $r=11$ remains. Hence $X\simeq\PP^{10}$.

\section{The index bound \texorpdfstring{$r\ge n-3$}{r >=  n-3}}\label{sec:high}

\subsection{The possible dimensions}
Put $c=n+1-r$, an even nonnegative integer. Substituting $n=r+c-1$ gives
\[
 \frac{n(n+1)^2}{2}-\frac{c^2(c-1)}2
 =\frac r2\bigl(r^2+(3c-1)r+3c^2-2c\bigr).
\]
The bracket is even: $r^2+(3c-1)r$ is even because $3c-1$ is odd,
and $3c^2-2c$ is even. Equation \eqref{eq:rcn} therefore implies
\begin{equation}\label{eq:fixedc}
 r\mid\frac{c^2(c-1)}2.
\end{equation}
For $r\ge n-3$, parity leaves $r=n+1,n-1,n-3$.
The first gives projective space.
For $r=n-1$, $c=2$ and \eqref{eq:fixedc} gives $r\mid2$,
hence $n=2,3$, covered by the established low-dimensional result.
For $r=n-3$, $c=4$ and $r\mid24$, so the full list is
\begin{equation}\label{eq:highpairs}
 (n,r)=(4,1),(5,2),(6,3),(7,4),(9,6),(11,8),(15,12),(27,24).
\end{equation}
The first three are again covered by $n\le6$.
We first exclude dimensions $7,9,11$ by Riemann--Roch, then treat
$15,27$ using the geometry of a curve section.

\subsection{Dimensions seven, nine, and eleven}\label{sec:higharith}
Let $m=h^0(L)$, with $0\le m\le n$, and $r=n-3$.
The Kodaira roots divide $P(t)$:
\[
 P(t)=\frac{\prod_{j=1}^{r-1}(t+j)}{n!}
       \bigl(z^2+\alpha z+\beta\bigr),\qquad z=t(t+r).
\]
To justify the shape, the quotient has degree four. Under $t\mapsto-r-t$
both $P$ and the product of roots pick up the same sign because
$n-(r-1)=4$ is even. The quotient is invariant, so it is a quadratic
polynomial in $z$; its leading coefficient is one.
The conditions $P(0)=1$, $P(1)=m$ give
\begin{equation}\label{eq:highP}
 \beta=r(r+1)(r+2)(r+3),\qquad
 \alpha=(r+2)(r+3)(m-r)-(r+1).
\end{equation}
For example, the product at $t=1$ is $r!$ and $z=r+1$;
substitution verifies the second formula directly.

Recover all Todd coefficients by \eqref{eq:recoverTodd}, form
$e^{-rx/2}\td$, and take its logarithm. Dividing the coefficient of $x^{2j}$
by $-B_{2j}/(2j(2j)!)$ determines $s_{2j}$.
Set $s_1=r$ and temporarily leave $s_3,s_5,\ldots,s_n$ unknown.

The coefficient of $s_{n-2}$ in $c_{n-1}$ is $r/(n-2)$.
Indeed, the generating identity
\[
 c(t)=\exp\left(\sum_{j\ge1}(-1)^{j-1}s_jt^j/j\right)
\]
gives $\partial c_{n-1}/\partial s_{n-2}=c_1/(n-2)$ in these odd dimensions.
Thus $c_{n-1}=n(n+1)^2/(2r)$ determines $s_{n-2}$ uniquely.
Next, $c_n=n+1$ determines $s_n$ with coefficient $1/n$.
The free odd variables left are
\[
 \begin{array}{c|c}
 n&\text{remaining variables}\\ \hline
 7&s_3\\9&s_3,s_5\\11&s_3,s_5,s_7.
 \end{array}
\]
Use \eqref{eq:exterior}--\eqref{eq:HRRexterior} for
$p=1,\ldots,(n-5)/2$. This gives one, two, or three polynomial equations.

We compute a Gr\"obner basis over $\Q$, ordering the remaining
variables lexicographically from the largest subscript down to $s_3$.
The basis contains
a nonzero polynomial solely in $s_3$.
After clearing denominators and a coefficient gcd, call it $E_{n,m}(s_3)$.
In each case, division by this Gr\"obner basis gives remainder zero.
Thus $E_{n,m}$ belongs to the rational ideal of the original equations
and vanishes at every common solution.
Consequently, the absence of integer roots of $E_{n,m}$ suffices to
exclude the original system.

For each $m$ the following table gives a prime modulo which $E_{n,m}$
has no root:
\[
\begin{array}{c|c|c|l}
 n&r&\deg E_{n,m}&\text{rootless primes, in order }m=0,\ldots,n\\ \hline
7&4&2&7,5,11,7,13,7,5,7\\
9&6&4&11,7,7,19,11,17,7,11,7,7\\
11&8&9&13,7,23,13,7,7,11,29,7,17,11,7.
\end{array}
\]
For instance,
\[
\begin{aligned}
 E_{9,0}(a)={}&13a^4-1496718a^3+541796355a^2\\
             &+87999724113156a+427746860206005240
\end{aligned}
\]
has no root modulo eleven.
For each of the $30$ polynomials, evaluation at every residue modulo
the listed prime gives a nonzero value. This excludes all three pairs.
The derivations, ideal-membership checks, and residue calculations
are performed by the verification program; see Section \ref{sec:python}.

\subsection{A Bernoulli-number congruence}
For every case $r=n-3$, put $\delta=n+1-h^0(L)$.
The reciprocal Hilbert numerator has degree four, and is forced to be
\begin{equation}\label{eq:hfour}
 h(z)=1-\delta z+(2\delta-1)z^2-\delta z^3+z^4.
\end{equation}
Its middle coefficient follows by evaluating $h(1)=1$.
Thus
\[
 B_\delta(x)=e^{2x}h(e^{-x})
 =2\cosh(2x)-2\delta\cosh x+2\delta-1.
\]
In $\log\Ahat=(n+1)\log q+\log B_\delta$, comparison of degree $2j$ gives
\begin{equation}\label{eq:powerBernoulli}
 s_{2j}=n+1-\frac{2j(2j)!}{B_{2j}}[x^{2j}]\log B_\delta(x).
\end{equation}
We apply this identity with $2j=12$, which is at most $n$ for the two
remaining dimensions $15$ and $27$.

Let $F(\delta)=12![x^{12}]\log B_\delta(x)$.
The recurrence \eqref{eq:logrec} gives
\begin{align*}
F(\delta)=-2\bigl(&39916800\delta^6-938044800\delta^5
                     +8944523280\delta^4\\
 &-44218331520\delta^3+119360516286\delta^2\\
 &-166654367087\delta+94009610240\bigr).
\end{align*}
Since $B_{12}=-691/2730$, equation \eqref{eq:powerBernoulli} says
\[
 s_{12}=n+1+\frac{32760}{691}F(\delta).
\]
The integer $32760$ is coprime to $691$. Integrality of $s_{12}$
therefore forces $691\mid F(\delta)$.
Modulo $691$ the polynomial factors as
\[
 F(\delta)\equiv
 -297\delta(\delta-1)(\delta-2)
 (\delta^3+325\delta^2-326\delta+237).
\]
Evaluation at the $691$ residues shows that the cubic has no root
in $\mathbf F_{691}$. Hence
$\delta\equiv0,1,2\pmod{691}$.
But $0\le\delta\le n+1\le28$ in \eqref{eq:highpairs}, so
\begin{equation}\label{eq:deltasmall}
 \delta\in\{0,1,2\}.
\end{equation}
The case $\delta=0$ gives $h^0(L)=n+1$ and is already impossible for
$r=n-3$ by Lemma \ref{lem:equality}. We treat $\delta=1,2$ by passing to a curve section.

\subsection{A curve section of genus two}
For $\delta=1,2$, we have $m=h^0(L)=n,n-1$.
Lemma \ref{lem:base} permits $n-1$ general independent sections
$s_1,\ldots,s_{n-1}$ to cut a proper complete intersection curve $C$.
Its intersection cycle has $L$-degree $L^n=1$.
Every irreducible component has positive integral degree, so there is a
single component of multiplicity one. A complete intersection in the
smooth $X$ is Cohen--Macaulay and has no embedded associated primes.
Since it is generically reduced, it is reduced everywhere: a nonzero
nilpotent submodule would have an associated prime, necessarily minimal,
contradicting generic reducedness. Thus $C$ is integral and Gorenstein.

Put $A=L|_C$. Adjunction for a local complete intersection gives
\[
 \omega_C=(K_X+(n-1)L)|_C=A^2,\qquad \deg A=1.
\]
Therefore $\deg\omega_C=2$ and $p_a(C)=2$.
The following argument applies to this integral Gorenstein curve even
when it is singular.

To identify the section ring of the curve, let $R=\bigoplus_{k\ge0}H^0(X,L^k)$.
For every integer $k$ and $0<i<n$, $H^i(X,L^k)=0$:
Kodaira vanishing handles $k>-r$, and Serre duality reduces $k\le-r$
to a nonnegative twist. The usual graded local-cohomology description
of a full section ring then makes $R$ Cohen--Macaulay of dimension $n+1$.
The $n-1$ sections have height $n-1$, so form a regular sequence.
The quotient $\bar R=R/(s_1,\ldots,s_{n-1})$ has dimension and depth two.
The exact comparison sequence with sections on $\Proj\bar R$ has kernel
$H^0_{R_+}(\bar R)$ and cokernel $H^1_{R_+}(\bar R)$, both zero by that
depth. Hence
\begin{equation}\label{eq:curve-ring}
 \bar R=\bigoplus_{k\ge0}H^0(C,A^k).
\end{equation}
Its degree-one dimension is
\[
 h^0(C,A)=m-(n-1)=
 \begin{cases}1&\delta=1,\\0&\delta=2.\end{cases}
\]

\subsection{Section rings of genus-two curves}\label{sec:curvelemma}
The presentations in Lemma \ref{lem:halfring} are those of
\cite[Proposition 3.2]{fpr}. We include a proof for integral Gorenstein
curves.

\begin{lemma}\label{lem:canonical}
For an integral Gorenstein curve $C$ of arithmetic genus two, $\omega_C$
is globally generated. Its two sections define a finite flat map
$f:C\to\PP^1$ of degree two, with $f^*\OO(1)=\omega_C$.
\end{lemma}
\begin{proof}
Suppose a point $p$ were a base point, and let $I_p$ be its ideal.
Then $h^0(\omega_C I_p)=2$ and $\chi(\omega_C I_p)=0$.
Serre duality for the torsion-free sheaf $\omega_C I_p$ gives
$h^0(\mathcal H om(I_p,\OO_C))=2$.

If $p$ is smooth, the latter sheaf is $\OO_C(p)$.
A nonconstant function with only a simple pole at $p$ defines a
degree-one morphism $C\to\PP^1$, necessarily an isomorphism,
contradicting $p_a(C)=2$.
If $p$ is singular, put $J=\mathcal H om(I_p,\OO_C)$.
Locally at $p$ it is the fractional ideal
$\{g:g\mathfrak m_p\subset\OO_{C,p}\}$.
In fact $g\mathfrak m_p\subset\mathfrak m_p$: otherwise $ga$ is a unit
for some $a\in\mathfrak m_p$, and every $b\in\mathfrak m_p$ equals
$a(gb)/(ga)$, making the maximal ideal principal and the point regular.
Thus $J=\mathcal E nd(I_p)$, a finite sheaf of algebras.
Its relative spectrum is an integral projective curve finite birational
over $C$. Such a curve has only constant global functions, contradicting
$h^0(J)=2$. There is no base point.

The canonical map is nonconstant since $\deg\omega_C=2>0$.
It is finite, has degree two, and pulls $\OO(1)$ back to $\omega_C$.
Its direct image of $\OO_C$ is torsion-free over the smooth curve $\PP^1$,
and hence locally free. This proves flatness.
\end{proof}

\begin{lemma}\label{lem:halfring}
Let $A$ have degree one and $A^2\simeq\omega_C$.
If $h^0(A)=1$, then $R(C,A)$ has generators of degrees $1,2,5$
and one relation of degree ten.
If $h^0(A)=0$, it has generators of degrees $2,2,3,3$
and two relations of degree six.
\end{lemma}
\begin{proof}
First suppose $h^0(A)=1$ and choose its section $x_1$.
Its zero scheme has length one and is an effective Cartier divisor.
The point $p$ is smooth: its local maximal ideal is generated by
the local equation of that divisor.
Thus $A=\OO_C(p)$.
Riemann--Roch and $A^2=\omega_C$ give
\[
 h^0(A^0)=1,\quad h^0(A)=1,\quad h^0(A^2)=2,\quad
 h^0(A^k)=k-1\quad(k\ge3).
\]
The pole orders at $p$ of functions regular off $p$ are therefore
$0,2,4,5,6,7,\ldots$; the missing positive orders are $1,3$.
Choose functions $y,z$ of exact pole orders $2,5$.
Every allowed pole order is $2b+5\epsilon$ with $\epsilon=0,1$.
Subtract a suitable multiple of $y^bz^\epsilon$ to reduce the pole order,
and iterate. This proves generation by $y,z$ for the affine function
ring off $p$, and by the homogeneous generators $x_1,y,z$ of degrees
$1,2,5$ for the section ring.
The square $z^2$ has pole order ten; reducing it gives a homogeneous
relation of degree ten with nonzero $y^5$ coefficient.
After using that relation to reduce the exponent of $z$ to $0$ or $1$,
the remaining monomials in any fixed degree have distinct pole orders.
They are linearly independent and form the full space counted above.
Thus there are no further relations.

Now suppose $h^0(A)=0$, and use the double cover in Lemma \ref{lem:canonical}.
The trace decomposition and Euler characteristics give
\[
 f_*\OO_C=\OO\oplus\OO(-3),\qquad
 f_*A=\OO(-1)\oplus\OO(-1).
\]
For the second equality, a rank-two vector bundle on $\PP^1$ splits as
a sum of line bundles. Its Euler characteristic is $\chi(A)=0$,
so its degree is $-2$; absence of sections forces both summands to have
negative degree, hence both are $-1$.

Let $S=\C[y_1,y_2]$, with both variables of degree two, corresponding to
$H^0(A^2)$. The projection formula gives, as graded $S$-modules,
\[
 R(C,A)_{\rm even}=S\oplus S w,\quad \deg w=6,\qquad
 R(C,A)_{\rm odd}=S z_1\oplus S z_2,\quad\deg z_i=3.
\]
Here $f_*A^3=\OO^2$ supplies $z_1,z_2$.
Finite-flat duality gives a perfect pairing
\[
 f_*A^3\otimes f_*A^3\longrightarrow
 f_*A^6=f_*\omega_{C/\PP^1}\longrightarrow\OO_{\PP^1}.
\]
The equality $A^6=\omega_{C/\PP^1}$ follows from
$\omega_C=A^2$ and $f^*\omega_{\PP^1}^{-1}=A^4$.
This pairing is symmetric. It is given by a nondegenerate constant
$2\times2$ matrix on $\OO^2$.
Over $\C$ choose an isotropic basis: the self-pairings of $z_1,z_2$
vanish and their mutual pairing is nonzero.
The trace kills the $\OO(3)$ summand of $f_*A^6=\OO(3)\oplus\OO$,
so $z_1^2$ and $z_2^2$ are cubic polynomials in $y_1,y_2$,
while $z_1z_2$ has nonzero $w$ coefficient.
It follows that
\[
 R(C,A)=\C[y_1,y_2,z_1,z_2]/
        (z_1^2-f_3(y_1,y_2),\,z_2^2-g_3(y_1,y_2)).
\]
To check that the displayed relations are complete, their quotient is free
over $S$ with basis $1,z_1,z_2,z_1z_2$.
These map to a basis of the already computed free $S$-module $R(C,A)$.
Thus the map is an isomorphism, proving the claim.
\end{proof}

\subsection{The section ring of the manifold}
By \eqref{eq:curve-ring}, the curve generators lift to $R$.
Together with $s_1,\ldots,s_{n-1}$ they generate $R$: reducing a homogeneous
element modulo the regular sequence and then subtracting its expression
reduces the remaining degrees, so induction proves surjectivity.

Relations lift as well. For one regular element $s$, a polynomial relation
$g$ in the quotient evaluates in $R$ to $s\,b$.
Express $b$ in the lifted generators and replace $g$ by $g-sB$;
this is a relation in $R$ reducing to $g$.
Repeat through the regular sequence.
To see that these generate all relations, observe that modulo the new degree-one variables the
lifted equations form the given complete intersection. Hence, together
with those variables, they form a regular sequence in the polynomial ring.
Its Hilbert series is the curve Hilbert series divided by $(1-z)^{n-1}$,
which is also the Hilbert series of $R$ by regularity of the original
sections. The surjection is therefore an isomorphism in each degree.

Consequently $X=\Proj R$ has one of the presentations
\begin{equation}\label{eq:weighted}
 X_{10}\subset\PP(1^n,2,5),\qquad
 X_{6,6}\subset\PP(1^{n-1},2,2,3,3).
\end{equation}
The notation $1^n$ means $n$ coordinates of weight one.
The first ambient space has dimension $n+1$ and one equation;
the second has dimension $n+2$ and two equations.
The equality $X=\Proj R$ follows for every ample line bundle by taking a
sufficiently high very ample Veronese subring, which does not change Proj.

\subsection{The Euler characteristic}
At a point of $X$, consider the weights of the homogeneous generators
which are nonzero there. A stabilizer of the weighted scalar action has
order equal to the gcd of those weights.
For all sufficiently large integers $k$, the line bundle $L^k$ is generated
by sections. Hence at that point, for every such $k$, some degree-$k$
polynomial in the generators is nonzero. A nonzero monomial uses only
generators nonzero there, so their gcd divides every sufficiently large
integer. Their gcd is therefore one.
Thus $X$ lies in the free-action locus of the weighted ambient space.

On this smooth locus the weighted Euler sequence is
\[
 0\longrightarrow\OO\longrightarrow\bigoplus_i\OO(w_i)
 \longrightarrow T_{\rm amb}\longrightarrow0.
\]
The first map is the weighted radial vector field.
The normal bundle of a smooth complete intersection of degrees $d_j$
is $\bigoplus_j L^{d_j}$. Restricting the Euler sequence and using the
normal sequence yields
\[
 c(T_X)=\frac{\prod_i(1+w_i x)}{\prod_j(1+d_j x)}.
\]
For \eqref{eq:weighted} this becomes
\begin{equation}\label{eq:weightedchern}
 \frac{(1+x)^n(1+2x)(1+5x)}{1+10x},
 \qquad
 \frac{(1+x)^{n-1}(1+2x)^2(1+3x)^2}{(1+6x)^2}.
\end{equation}

We compute the top coefficients by a generating-function identity.
For any series $F(x)$,
\[
 \sum_{n\ge0}\bigl([x^n](1+x)^nF(x)\bigr)t^n
 =\frac1{1-t}F\!\left(\frac{t}{1-t}\right).
\]
It follows by summing $\sum_{n\ge j}\binom njt^n=t^j/(1-t)^{j+1}$.
For the first expression in \eqref{eq:weightedchern}, subtracting the
generating series $\sum_{n\ge0}(n+1)t^n=(1-t)^{-2}$ gives
\[
 \frac{(1+t)(1+4t)}{(1-t)^2(1+9t)}-\frac1{(1-t)^2}
 =\frac25\left(\frac1{1+9t}-\frac1{1-t}\right).
\]
For the second it gives
\[
\begin{aligned}
 \frac{(1+t)^2(1+2t)^2}{(1-t)^2(1+5t)^2}-\frac1{(1-t)^2}
 ={}&\frac4{25}+\frac{26}{75(1+5t)}\\
 &+\frac4{25(1+5t)^2}-\frac2{3(1-t)}.
\end{aligned}
\]
Extracting any positive degree $n$ yields respectively
\[
 c_n-(n+1)=\frac25\bigl((-9)^n-1\bigr),\qquad
 c_n-(n+1)=\frac{(12n+38)(-5)^n-50}{75}.
\]
Both are negative when $n=15$ or $27$, contradicting $c_n=n+1$.
This excludes the two remaining pairs in \eqref{eq:highpairs}.
This proves Theorem \ref{thm:index}\textup{(i)}.

\section{The index bound in even dimensions}\label{sec:even}

\subsection{The case \texorpdfstring{$r=n-5$}{r=n-5}}\label{sec:even6}
Now $c=n+1-r=6$, so \eqref{eq:fixedc} gives $r\mid90$.
If $n$ is even, $r$ is odd. The odd divisors $1,3,5,9,15,45$
therefore give exactly
\[
 (n,r)=(6,1),(8,3),(10,5),(14,9),(20,15),(50,45).
\]
Let $m=h^0(L)$ and $k=h^0(2L)$, and put
\[
 A=m-n-1,\qquad B=k-(n+1)m+\frac{n(n+1)}2.
\]
Reciprocity and $h(1)=1$ give
\[
 h(z)=1+Az+Bz^2-(1+2A+2B)z^3+Bz^4+Az^5+z^6.
\]
Substitution in \eqref{eq:ND} gives
\begin{align*}
 N(t)={}&1+(16A+4B+33)t^2+(-16A-8B+27)t^4+(4B+3)t^6,\\
 D(t)={}&1+(64A+16B+137)t^2+(-64A-80B+837)t^4\\
 &+(-128A+160B+893)t^6+(128A-160B-893)t^8\\
 &+(64A+80B-837)t^{10}+(-64A-16B-137)t^{12}-t^{14}.
\end{align*}
Compute $a_{n/2}$ by \eqref{eq:division}; the required equation is
$a_{n/2}=1$.

For each $m=0,\ldots,n$, a prime $p$ is specified in
Appendix \ref{app:evenprimes}. For every $k=0,\ldots,p-1$ the computed
$a_{n/2}$ differs from $1$ modulo $p$.
That excludes \emph{every} integer $k$, with no upper bound required.
The numbers of specializations and the largest primes used are
\[
\begin{array}{r|r|r|r}
n&r&\text{values of }m&\text{largest prime}\\ \hline
6&1&7&7\\8&3&9&19\\10&5&11&13\\14&9&15&23\\
20&15&21&19\\50&45&51&61
\end{array}
\]
The $114$ pairs $(n,m)$ are therefore excluded. In particular,
this proves the exclusion of $n=10,r=5$ used in Section \ref{sec:C10}.

\subsection{The case \texorpdfstring{$r=n-7$}{r=n-7}}
Here $c=8$, and \eqref{eq:fixedc} gives $r\mid224$.
The only odd divisors of $224$ are $1,7$, so the only even-dimensional
pairs are $(8,1)$ and $(14,7)$. Section \ref{sec:C8} already excludes
the former.

For $(14,7)$ set $m=h^0(L)$, $k=h^0(2L)$, $\ell=h^0(3L)$.
The first four coefficients of $h(z)$ are
\[
 h_0=1,\quad h_1=m-15,\quad h_2=k-15m+105,\quad
 h_3=\ell-15k+105m-455.
\]
Set $h_4=1-2\sum_{i=0}^3h_i$ and reflect.
The integer equation obtained from
$[t^{14}]N(t)^2/D(t)=1$, after dividing its coefficient gcd,
is denoted $F_{14,7}(m,k,\ell)=0$. It has total degree seven and 120
monomials. It is obtained by \eqref{eq:ND} and \eqref{eq:division}, and its
coefficients are listed in
Appendix \ref{app:n14poly}.

The section bound and the equality lemma give
$0\le m\le14$ and $0\le k\le2^{14}+14=16398$.
For each pair $(m,k)$ we test the resulting polynomial in $\ell$
modulo successive primes, obtaining
\[
\begin{array}{r|r}
\text{primes tested through}&\text{remaining pairs}\\ \hline
\text{initial}&245985\\31&4299\\61&175\\97&5\\113&0.
\end{array}
\]
Thus $(14,7)$ is excluded as well.

Finally, if $n$ is even and $r\ge n-7$, parity gives just
$r=n+1,n-1,n-3,n-5,n-7$, omitting any nonpositive value.
The previous section treats the first three; this section excludes the last
two. This proves Theorem \ref{thm:index}\textup{(ii)}.
The proofs of $C_7$ through $C_{10}$, together with the established
$n\le6$ theorem, prove Theorem \ref{thm:dimension}.

\section{Computational verification}\label{sec:python}

\subsection{The verification program}

The verification program
\nolinkurl{verify_fujita_further_extensions.py} uses SymPy for
polynomial algebra over $\Q$ and Python integers for the finite
congruence calculations. It constructs the Hilbert polynomials,
characteristic classes, and eliminants from the formulas in the
preceding sections. The parameter bounds and moduli are specified in
the source and in the tables of this article; no input data files
are required.

The calculation has two parts. First, the program derives the
polynomial equations using Newton's identities, Riemann--Roch, and
the signature formula. It then excludes the permitted integer
solutions by the discriminant and congruence tests. A successful run
verifies these algebraic identities and finite calculations.
The geometric arguments on which they depend are given in the
preceding sections.

The principal functions correspond to the text as follows.
\begin{center}
\small
\begin{tabular}{@{}>{\raggedright\arraybackslash}p{0.18\textwidth}>{\raggedright\arraybackslash}p{0.44\textwidth}>{\raggedright\arraybackslash}p{0.31\textwidth}@{}}
\toprule
Calculation&Functions&Assertions checked\\
\midrule
Dimension seven&
\nolinkurl{c7_symbolic_certificate}\par
\nolinkurl{c7_finite_certificate}&
Chern identities, both eliminations, and discriminants.\\[4pt]
Dimension eight&
\nolinkurl{c8_symbolic_certificate}\par
\nolinkurl{c8_arithmetic_certificate}&
Hilbert polynomials, signature quartics, and congruences.\\[4pt]
Dimension nine&
\nolinkurl{c9_eliminant}\par
\nolinkurl{sieve_three_parameters}&
The resultant quartic and all bounded triples.\\[4pt]
Dimensions ten and fourteen&
\nolinkurl{signature_polynomial}\par
\nolinkurl{compact_c10_polynomial}\par
\nolinkurl{check_signature_reconstruction}\par
\nolinkurl{sieve_three_parameters}\par
\nolinkurl{small_signature_sieve}&
Signature polynomials and the corresponding congruence tests.\\[4pt]
$r=n-3$&
\nolinkurl{high_index_arithmetic}\par
\nolinkurl{high_bernoulli_certificate}\par
\nolinkurl{high_dimension_and_euler_certificates}&
The $30$ eliminants, the congruence modulo $691$, and the Euler-characteristic formulas.\\[4pt]
$r=n-5$, $n$ even&
\nolinkurl{coindex6_even}&
The signature recurrence for all $114$ specializations.\\
\bottomrule
\end{tabular}
\end{center}

For the symbolic calculations, the program constructs the power series
by \eqref{eq:logrec} and \eqref{eq:exprec}, using rational coefficients.
It solves linear equations, forms the dimension-nine resultant, and
computes the Gr\"obner bases in Section \ref{sec:higharith} over $\Q$.
Denominators and common coefficient factors are removed before an
integer polynomial is passed to a congruence test.

A polynomial in one variable is evaluated modulo $p$ by Horner's rule,
at every residue $0,\ldots,p-1$, until a root is found or the residues
are exhausted. A polynomial which becomes identically zero modulo $p$
is retained, as it must be. The procedure also accommodates a drop in
degree after reduction. In dimension seven, the alternative
discriminant test uses the integer square root.

\subsection{Representation of the candidate sets}

The two largest initial sets have $102797460$ and $672398100$ triples.
The program represents a set of third coordinates by the binary digits
of a nonnegative integer. For fixed $(m,k)$, write
\[
 B_{m,k}=\sum_{\ell\ \mathrm{allowed}}2^\ell.
\]
Initially $0\le\ell<L$ and $B_{m,k}=2^L-1$. This represents all
possible third coordinates for the given pair.

Fix a prime $p$. Whether the specialized polynomial has a root modulo
$p$ depends only on the parameter residues. For each pair
$(m,k\bmod p)$, form
\[
 R=\sum_{\substack{0\le a<p\\
 \text{the specialization at }\ell=a\text{ has a root modulo }p}}2^a.
\]
If $q=\lceil L/p\rceil$, the integer
\[
 M=R\left(1+2^p+\cdots+2^{(q-1)p}\right)
       \mathbin{\&}(2^L-1)
  =R\,\frac{2^{pq}-1}{2^p-1}\mathbin{\&}(2^L-1)
\]
has digit $\ell$ equal to digit $\ell\bmod p$ of $R$. Here $\&$
denotes bitwise intersection. The repeated blocks have disjoint
supports, so no carries occur in the multiplication. Replacing
$B_{m,k}$ by $B_{m,k}\mathbin{\&}M$ therefore applies the congruence
condition to every permitted third coordinate.

The number of remaining triples is the sum of the numbers of nonzero
binary digits, computed by \nolinkurl{int.bit_count()}. Once at most
$500000$ triples remain, the program converts the set digits to a list
of triples and continues by direct evaluation. The length of this list
is checked against the preceding count. Thus both representations
describe the same finite set at every stage.

The periodic repetition formula is implemented in
\nolinkurl{repeat_residues}. The program checks it against direct
enumeration for every subset of residues at moduli $1,\ldots,7$ and
lengths $1,\ldots,19$. It also compares the two representations on a
small range using the dimension-nine polynomial itself.

At each prime, evaluations for a repeated tuple of parameter residues
are stored and reused. This is valid because evaluation of an integer
polynomial commutes with reduction. For dimension nine, the
substitutions \eqref{eq:C9powers} also have integer coefficients and
therefore commute with reduction. The stored values are discarded when
the prime changes.

\subsection{Reproduction of the calculations}

The source file requires Python 3.10 or later and SymPy. Running
\nolinkurl{python verify_fujita_further_extensions.py} derives the
equations and executes all the finite checks. It prints the successive
counts and ends with \texttt{ALL PYTHON CHECKS PASSED.} if every
assertion succeeds. The tests use Python assertions, so the program is
run without the optimization flag \texttt{-O}.

The option \texttt{--export-json audit.json} additionally writes the
derived coefficient arrays, eliminants, residue values, and sieve
counts. This file is an output of the calculation. The high-index
eliminants have coefficient lists in descending powers. Other
polynomials are recorded by their variables, monomial exponents, and
integer coefficients; when present, the additional
\nolinkurl{coefficients_latex} list runs in increasing powers of the
specified root variable. The program also checks its
characteristic-class recurrences on $\PP^7$, $\PP^9$, and $\PP^{11}$,
and verifies the signature polynomials by recovering the Todd class
from the Hilbert polynomial.

The repository contains the TeX source, this article, the Python file,
and instructions for compiling and running them. The polynomial
appendices below allow the equations and congruence tables to be read
independently of the source code.

\appendix

\clearpage
\section{The eliminant in dimension nine}\label{app:c9poly}
We use the notation of Section \ref{sec:C9}:
$u=s_2,v=s_4,w=s_6,a=s_3,b=s_5$.
The two primitive equations obtained after eliminating $s_7,s_8,s_9$
are $E_i=L_i b+C_i=0$, with the following coefficients.
\begin{align*}
L_1&={}-8640 u^{2} - 6336 a u\\
&\quad{} + 410112 u - 17280 v\\
&\quad{} + 103680 a - 1078272
\end{align*}
\begin{align*}
C_1&={}-14685 u^{5} + 324060 u^{4}\\
&\quad{} - 31680 u^{3} v - 240 a u^{3}\\
&\quad{} - 2326560 u^{3} + 590400 u^{2} v\\
&\quad{} - 34320 u^{2} w + 2640 a^{2} u^{2}\\
&\quad{} - 160320 a u^{2} + 6453120 u^{2}\\
&\quad{} - 9900 u v^{2} + 30240 a u v\\
&\quad{} - 3127680 u v + 215040 u w\\
&\quad{} - 53760 a^{2} u + 890112 a u\\
&\quad{} + 114974961216 u + 162000 v^{2}\\
&\quad{} - 362880 a v + 3663360 v\\
&\quad{} - 1920 a w - 295680 w\\
&\quad{} + 640 a^{3} + 157440 a^{2}\\
&\quad{} - 844800 a - 844738332416
\end{align*}
\begin{align*}
L_2&={}-43200 u^{2} + 113472 a u\\
&\quad{} + 889344 u - 86400 v\\
&\quad{} - 1610496 a + 2350080
\end{align*}
\begin{align*}
C_2&={}262995 u^{5} - 5823780 u^{4}\\
&\quad{} + 567360 u^{3} v + 59280 a u^{3}\\
&\quad{} + 41549280 u^{3} - 8781120 u^{2} v\\
&\quad{} + 614640 u^{2} w - 47280 a^{2} u^{2}\\
&\quad{} - 559680 a u^{2} - 112402560 u^{2}\\
&\quad{} + 177300 u v^{2} - 211680 a u v\\
&\quad{} + 39035520 u v - 5859840 u w\\
&\quad{} + 860160 a^{2} u - 1807104 a u\\
&\quad{} - 2059103882688 u - 702000 v^{2}\\
&\quad{} - 362880 a v - 39744000 v\\
&\quad{} + 474240 a w + 9811200 w\\
&\quad{} - 158080 a^{3} + 1109760 a^{2}\\
&\quad{} + 678912 a + 13261343095040
\end{align*}

Their resultant satisfies
\[
 \Res_b(E_1,E_2)=241920F(u,v,w,a).
\]
The coefficients of $F=\sum_{i=0}^4A_i(u,v,w)a^i$ are
\begin{align*}
A_4&={}3840 u - 63488
\end{align*}
\begin{align*}
A_3&={}5760 u^{2} - 362496 u\\
&\quad{} + 11520 v + 2222080
\end{align*}
\begin{align*}
A_2&={}720 u^{4} - 16512 u^{3}\\
&\quad{} - 4320 u^{2} v + 152064 u^{2}\\
&\quad{} + 209664 u v - 11520 u w\\
&\quad{} + 3520512 u - 2594304 v\\
&\quad{} + 190464 w - 11808768
\end{align*}
\begin{align*}
A_1&={}13320 u^{5} - 242784 u^{4}\\
&\quad{} - 6048 u^{3} v + 1837824 u^{3}\\
&\quad{} + 72576 u^{2} v + 70272 u^{2} w\\
&\quad{} - 7492608 u^{2} - 21600 u v^{2}\\
&\quad{} - 2953728 u v - 387072 u w\\
&\quad{} - 68163073536 u + 673920 v^{2}\\
&\quad{} - 34560 v w + 12146688 v\\
&\quad{} + 141312 w + 59894179840
\end{align*}
\begin{align*}
A_0&={}-12015 u^{7} + 765684 u^{6}\\
&\quad{} - 49950 u^{5} v - 13992912 u^{5}\\
&\quad{} + 2029032 u^{4} v - 28080 u^{4} w\\
&\quad{} + 106964928 u^{4} - 59940 u^{3} v^{2}\\
&\quad{} - 25028928 u^{3} v + 1415808 u^{3} w\\
&\quad{} + 93693805248 u^{3} + 1229040 u^{2} v^{2}\\
&\quad{} - 56160 u^{2} v w + 123482880 u^{2} v\\
&\quad{} - 13533696 u^{2} w - 4537368421632 u^{2}\\
&\quad{} - 16200 u v^{3} - 6384960 u v^{2}\\
&\quad{} + 495360 u v w + 187916888448 u v\\
&\quad{} + 41748480 u w + 33647373560832 u\\
&\quad{} + 108000 v^{3} + 5702400 v^{2}\\
&\quad{} - 806400 v w - 1248789496320 v\\
&\quad{} - 40857600 w - 50901671137280
\end{align*}

Each displayed equation is a polynomial over the integers. The resultant
identity may also be verified simply by expanding $L_1C_2-L_2C_1$.

\clearpage
\section{The signature polynomial in dimension fourteen}\label{app:n14poly}
For $n=14,r=7$, write
$F_{14,7}(m,k,\ell)=\sum_{i=0}^7B_i(m,k)\ell^i$.
The variables are the section counts $m=h^0(L)$, $k=h^0(2L)$,
$\ell=h^0(3L)$, not coefficients of the Hilbert numerator.
The following coefficients, together with that convention, specify the
entire polynomial used in the sieve.
\begin{align*}
B_7&={}-18874368
\end{align*}
\begin{align*}
B_6&={}-7134511104 m + 1453326336 k\\
&\quad{} + 20362690560
\end{align*}
\begin{align*}
B_5&={}-1155790798848 m^{2} + 470877732864 k m\\
&\quad{} + 6596407590912 m - 47959769088 k^{2}\\
&\quad{} - 1343814893568 k - 9410722832384
\end{align*}
\begin{align*}
B_4&={}-104021171896320 m^{3} + 63568493936640 k m^{2}\\
&\quad{} + 890365964451840 m^{2} - 12949137653760 k^{2} m\\
&\quad{} - 362769292984320 k m - 2540044282281984 m\\
&\quad{} + 879262433280 k^{3} + 36951535779840 k^{2}\\
&\quad{} + 517495201488896 k + 2415126819309568
\end{align*}
\begin{align*}
B_3&={}-5617143282401280 m^{4} + 4576931563438080 k m^{3}\\
&\quad{} + 64095617097400320 m^{3} - 1398506866606080 k^{2} m^{2}\\
&\quad{} - 39172524988170240 k m^{2} - 274232905622028288 m^{2}\\
&\quad{} + 189920685588480 k^{3} m + 7980195706306560 k^{2} m\\
&\quad{} + 111741528791384064 k m + 521405253005328384 m\\
&\quad{} - 9671886766080 k^{4} - 541906375802880 k^{3}\\
&\quad{} - 11382814465998848 k^{2} - 106236446111965184 k\\
&\quad{} - 371714603757760192
\end{align*}
\begin{align*}
B_2&={}-181995442349801472 m^{5} + 185365728319242240 k m^{4}\\
&\quad{} + 2595437832547860480 m^{4} - 75519370796728320 k^{2} m^{3}\\
&\quad{} - 2114962182037831680 k m^{3} - 14803616548748132352 m^{3}\\
&\quad{} + 15383575532666880 k^{3} m^{2} + 646287634417582080 k^{2} m^{2}\\
&\quad{} + 9048032227521921024 k m^{2} + 42212600683249065984 m^{2}\\
&\quad{} - 1566845656104960 k^{4} m - 87774136636538880 k^{3} m\\
&\quad{} - 1843398338061828096 k^{2} m - 17201655546463420416 k m\\
&\quad{} - 60177407870074533248 m + 63834452656128 k^{5}\\
&\quad{} + 4470319371386880 k^{4} + 125188081996054528 k^{3}\\
&\quad{} + 1752420756180602880 k^{2} + 12262096027360696000 k\\
&\quad{} + 34310884872797383402
\end{align*}
\begin{align*}
B_1&={}-3275917962296426496 m^{6} + 4003899731695632384 k m^{5}\\
&\quad{} + 56052068529261772800 m^{5} - 2039023011511664640 k^{2} m^{4}\\
&\quad{} - 57094416397290700800 k m^{4} - 399563744123097317376 m^{4}\\
&\quad{} + 553808719176007680 k^{3} m^{3} + 23262458998475980800 k^{2} m^{3}\\
&\quad{} + 325620044334329757696 k m^{3} + 1518889838325936979968 m^{3}\\
&\quad{} - 84609665429667840 k^{4} m^{2} - 4739009781222604800 k^{3} m^{2}\\
&\quad{} - 99510166256361209856 k^{2} m^{2} - 928422555999352602624 k m^{2}\\
&\quad{} - 3247402009645454876928 m^{2} + 6894120886861824 k^{5} m\\
&\quad{} + 482713662770380800 k^{4} m + 13515784244957085696 k^{3} m\\
&\quad{} + 189166326551280304128 k^{2} m + 1323418788249959729408 k m\\
&\quad{} + 3702470602277503658824 m - 234059659739136 k^{6}\\
&\quad{} - 19667609026560000 k^{5} - 688408640537829376 k^{4}\\
&\quad{} - 12847561883227873280 k^{3} - 134833727662115152448 k^{2}\\
&\quad{} - 754494298636214961516 k - 1758664553769502206697
\end{align*}
\begin{align*}
B_0&={}-25271367137715290112 m^{7} + 36035097585260691456 k m^{6}\\
&\quad{} + 504384118879407833088 m^{6} - 22021448524325978112 k^{2} m^{5}\\
&\quad{} - 616516421899247419392 k m^{5} - 4313842577830193135616 m^{5}\\
&\quad{} + 7476417708876103680 k^{3} m^{4} + 313990602631906590720 k^{2} m^{4}\\
&\quad{} + 4394397830215019397120 k m^{4} + 20494705844786211127296 m^{4}\\
&\quad{} - 1522973977734021120 k^{4} m^{3} - 85287891313297981440 k^{3} m^{3}\\
&\quad{} - 1790582921129257205760 k^{2} m^{3} - 16703206200200828583936 k m^{3}\\
&\quad{} - 58414055525899661578752 m^{3} + 186141263945269248 k^{5} m^{2}\\
&\quad{} + 13031086502636421120 k^{4} m^{2} + 364803926672740515840 k^{3} m^{2}\\
&\quad{} + 5104923740232864915456 k^{2} m^{2} + 35708356992589914385152 k m^{2}\\
&\quad{} + 99882939615801221882664 m^{2} - 12639221625913344 k^{6} m\\
&\quad{} - 1061873062887555072 k^{5} m - 37161614273488404480 k^{4} m\\
&\quad{} - 693419662686368464896 k^{3} m - 7276141194137848482176 k^{2} m\\
&\quad{} - 40708552503242224866488 k m - 94872385352327729726730 m\\
&\quad{} + 367808036732928 k^{7} + 36053990168199168 k^{6}\\
&\quad{} + 1514222256517341184 k^{5} + 35321107231167932416 k^{4}\\
&\quad{} + 494209522557826902080 k^{3} + 4147821048038964926314 k^{2}\\
&\quad{} + 19334692281325631968759 k + 38615228954566746978894
\end{align*}

\clearpage
\section{Congruence moduli for \texorpdfstring{$r=n-5$}{r=n-5}}
\label{app:evenprimes}
Each row says: for the indicated $n$ and each listed value of $m$, the
polynomial $a_{n/2}-1$ from the recurrence in Section \ref{sec:even6}
has no root $k$ modulo the indicated prime.
The lists for each $n$ partition $m=0,\ldots,n$.
\begin{longtable}{rrp{0.66\linewidth}}
\toprule
$n$&Prime&Values of $m$\\ \midrule
\endfirsthead
\toprule $n$&Prime&Values of $m$\\ \midrule\endhead
\bottomrule\endfoot
6&3&$1,2,4,5$\\
6&5&$0,3$\\
6&7&$6$\\
8&7&$2,3,5,6$\\
8&11&$0,1,4,8$\\
8&19&$7$\\
10&3&$0,1,3,4,6,7,9,10$\\
10&7&$2,5$\\
10&13&$8$\\
14&5&$1,2,6,7,11,12$\\
14&7&$4,5,8$\\
14&11&$0,3,14$\\
14&23&$9,10,13$\\
20&3&$1,4,7,10,13,16,19$\\
20&7&$3,17$\\
20&11&$2,12,15$\\
20&13&$5,6,18$\\
20&17&$8,9,11,14,20$\\
20&19&$0$\\
50&5&$1,6,11,16,21,26,31,36,41,46$\\
50&7&$4,5,8,12,15,18,19,22,25,29,32,33,39,40,43,47,50$\\
50&11&$0,2,10,13,24,27,35,38,44,49$\\
50&13&$34$\\
50&17&$7,28,30,45$\\
50&19&$37,48$\\
50&23&$14,23$\\
50&29&$20$\\
50&37&$42$\\
50&47&$3,17$\\
50&61&$9$\\
\end{longtable}

\clearpage
\section{Sieve tables}\label{app:sievetables}
The prime in a row is applied after all preceding primes in that column.
A dash means that prime was not used in that particular sieve.
The initial counts are $102797460$ for $C_9$ and $672398100$ for $C_{10}$.
\begin{longtable}{rrr}
\toprule Prime&Remaining triples for $C_9$&Remaining triples for $C_{10}$\\\midrule
\endfirsthead
\toprule Prime&Remaining triples for $C_9$&Remaining triples for $C_{10}$\\\midrule\endhead
\bottomrule\endfoot
3&--&448265745\\
5&--&267329529\\
7&--&181171361\\
11&--&119711405\\
13&71051845&80209438\\
17&47418292&51982239\\
19&31752841&33510036\\
23&20583262&21301530\\
29&13643896&13548080\\
31&8899360&8661101\\
37&5721568&5489690\\
41&3671815&3474806\\
43&2357661&2211620\\
47&1511609&1397055\\
53&964275&887441\\
59&611168&568511\\
61&390389&360233\\
67&248724&227836\\
71&157634&144692\\
73&99718&91881\\
79&63012&58870\\
83&40028&37415\\
89&25330&23681\\
97&16021&14926\\
101&10047&9501\\
103&6324&6032\\
107&3958&3869\\
109&2444&2403\\
113&1501&1511\\
127&938&946\\
131&579&614\\
137&357&409\\
139&217&254\\
149&136&148\\
151&87&90\\
157&57&63\\
163&39&41\\
167&25&28\\
173&14&18\\
179&10&12\\
181&3&7\\
191&3&3\\
193&2&3\\
197&2&2\\
199&1&1\\
211&1&1\\
223&0&1\\
227&--&1\\
229&--&1\\
233&--&1\\
239&--&0\\
\end{longtable}

For $n=14,r=7$, the complete pair counts are:
\begin{longtable}{rr}
\toprule Prime&Remaining pairs\\\midrule
\endfirsthead
\toprule Prime&Remaining pairs\\\midrule\endhead
\bottomrule\endfoot
2&245985\\
3&191325\\
5&160717\\
7&101680\\
11&63341\\
13&41844\\
17&28247\\
19&18531\\
23&10837\\
29&6812\\
31&4299\\
37&2733\\
41&1681\\
43&1069\\
47&659\\
53&415\\
59&249\\
61&175\\
67&117\\
71&68\\
73&40\\
79&24\\
83&13\\
89&10\\
97&5\\
101&4\\
103&2\\
107&2\\
109&1\\
113&0\\
\end{longtable}

\clearpage
\section{An eliminant modulo eleven}
The $n=9,r=6,m=0$ eliminant displayed in Section \ref{sec:higharith}
has the following values modulo eleven.
All entries are nonzero, so this specialization has no integer root.
\[
\begin{array}{c|rrrrrrrrrrr}
a&0&1&2&3&4&5&6&7&8&9&10\\\hline
E_{9,0}(a)\bmod11&
10&8&7&6&8&9&9&1&4&8&7\end{array}\]

The optional JSON export records every coefficient vector and every residue
list for all 30 high-index eliminants. In each coefficient vector the
coefficients are ordered from the highest power down to the constant term.
The accompanying program derives these polynomials from the Hilbert
polynomial.


\begin{thebibliography}{LW90}

\bibitem[Deb15]{debarre}
O.~Debarre,
\emph{Cohomological characterizations of the complex projective space},
arXiv:1512.04321, 2015.
\url{https://arxiv.org/abs/1512.04321}.

\bibitem[FPR17]{fpr}
M.~Franciosi, R.~Pardini, and S.~Rollenske,
\emph{Gorenstein stable surfaces with $K_X^2=1$ and $p_g>0$},
Math. Nachr. \textbf{290} (2017), no.~5--6, 794--814.
\url{https://arxiv.org/abs/1511.03238}.

\bibitem[Fuj80]{fujita}
T.~Fujita,
\emph{On topological characterizations of complex projective spaces and
affine linear spaces},
Proc. Japan Acad. Ser. A Math. Sci. \textbf{56} (1980), no.~5, 231--234.
\url{https://doi.org/10.3792/pjaa.56.231}.

\bibitem[KO73]{ko}
S.~Kobayashi and T.~Ochiai,
\emph{Characterizations of complex projective spaces and hyperquadrics},
J. Math. Kyoto Univ. \textbf{13} (1973), 31--47.

\bibitem[LP25]{lp}
P.~Li and T.~Peternell,
\emph{Compactification of homology cells, Fujita's conjectures and the
complex projective space},
arXiv:2502.01072v3, 2025.
\url{https://arxiv.org/abs/2502.01072v3}.

\bibitem[LW90]{lw}
A.~S.~Libgober and J.~W.~Wood,
\emph{Uniqueness of the complex structure on K\"ahler manifolds of certain
homotopy types},
J. Differential Geom. \textbf{32} (1990), no.~1, 139--154.
\url{https://doi.org/10.4310/jdg/1214445041}.

\end{thebibliography}

\end{document}
